\documentclass[10pt,a4paper]{article}

% Situación de aprendizaje de Diseño Curricular de las Matemáticas.
% Compilar con: latexmk -pdf Situacion_Aprendizaje_JGL_DCM.tex
% Las matrices (matrices_tabla.tex) se regeneran con: python3 matrices.py

% ---------- Partes opcionales ----------
% El reto y producto final y el desarrollo de las sesiones están en secciones/.
% Cada parte se activa o se oculta por separado con true o false.
\newif\ifproducto
\productotrue
\newif\ifsesiones
\sesionesfalse
\newcommand{\soloProducto}[1]{\ifproducto#1\fi}

% La justificación pedagógica de las relaciones (matrices_justificacion.tex) se conserva
% pero no se compila. Para mostrarla, cambia \justificacionfalse por \justificaciontrue.
\newif\ifjustificacion
\justificacionfalse
\newcommand{\soloJustificacion}[1]{\ifjustificacion#1\fi}

% ---------- Datos del documento ----------
\newcommand{\tituloSA}{Medir lo que no se puede alcanzar: ¿dónde caben más paneles solares en nuestro instituto?}
\newcommand{\autor}{Javier Gómez López}
\newcommand{\asignatura}{Diseño Curricular de las Matemáticas}
\newcommand{\master}{Máster Universitario en Formación del Profesorado -- Especialidad en Matemáticas}

% ---------- Paquetes ----------
\usepackage{float}
\usepackage[utf8]{inputenc}
\usepackage[T1]{fontenc}
\usepackage[spanish,es-noshorthands,es-tabla]{babel}
\usepackage[scaled=0.92]{helvet} % equivalente libre de la Arial del original
\renewcommand{\familydefault}{\sfdefault}
\usepackage[left=2cm,right=2cm,top=2.55cm,bottom=1.6cm,
            headheight=44pt,headsep=4pt,footskip=20pt]{geometry}
\usepackage{amsmath,amssymb}
\usepackage{xcolor}
\usepackage{graphicx}
\usepackage{fancyhdr}
\usepackage{lastpage}
\usepackage{etoolbox}
\usepackage{enumitem}
\usepackage{titlesec}
\usepackage{pdflscape}
\usepackage{tikz}
\usetikzlibrary{angles,quotes,arrows.meta}
\usepackage{tabularray}
\usepackage[most]{tcolorbox}
\usepackage{csquotes}
\usepackage[backend=bibtex,style=authoryear,maxcitenames=2,uniquename=false,uniquelist=false]{biblatex}
\addbibresource{bibliografia.bib}
\usepackage[hidelinks]{hyperref}

\setlength{\parindent}{0pt}
\setlength{\parskip}{5pt}
\setlist{nosep,leftmargin=1.2em}
\emergencystretch=1.5em

% ---------- Colores (medidos sobre la plantilla) ----------
\definecolor{urjc}{HTML}{C8102E}
\definecolor{rosa}{HTML}{FBE9EB}
\definecolor{borde}{HTML}{D0D0D0}
\definecolor{titular}{HTML}{3C3C3B}
\definecolor{pistagris}{HTML}{7F7F7F}

% ---------- Títulos ----------
\titleformat{\section}{\Large\bfseries\color{urjc}}{\thesection}{0.6em}{}[{\color{urjc}\titlerule[0.7pt]}]
\titleformat{\subsection}{\large\bfseries\color{titular}}{\thesubsection}{0.6em}{}
\titleformat{\subsubsection}{\normalsize\bfseries\color{urjc}}{}{0pt}{}
\titlespacing*{\section}{0pt}{14pt}{6pt}
\titlespacing*{\subsection}{0pt}{10pt}{3pt}
\titlespacing*{\subsubsection}{0pt}{8pt}{2pt}
\setcounter{secnumdepth}{2}
\setcounter{tocdepth}{2}
\makeatletter\renewcommand*\l@subsection{\@dottedtocline{2}{1.5em}{2.9em}}\makeatother

% ---------- Cabecera y pie ----------
\newcommand{\cabecera}{Máster en Formación del Profesorado · Diseño Curricular de las Matemáticas · 2026-27}

\pagestyle{fancy}
\fancyhf{}
\fancyhead[L]{\includegraphics[height=42pt]{logo_urjc.png}}
\fancyhead[R]{\footnotesize\color{pistagris}\cabecera}
\fancyfoot[R]{\footnotesize\color{pistagris}Página \thepage\ de \pageref{LastPage}}
\renewcommand{\headrule}{{\color{urjc}\hrule height 1pt}}
\renewcommand{\footrulewidth}{0pt}

% ---------- Elementos de texto ----------
% Etiqueta de un campo (Título:, Producto:...)
\newcommand{\campo}[1]{{\bfseries\color{urjc}#1}}
% Texto orientativo, en cursiva gris
\newcommand{\pista}[1]{{\itshape\small\color{pistagris}#1}}
% Lista con viñetas; dentro de las tablas hay que usar esta y no itemize
\newenvironment{lista}
  {\begin{minipage}[t]{\linewidth}\begin{itemize}[leftmargin=*,nosep]}
  {\end{itemize}\end{minipage}}
% Encabezado de los apartados a todo el ancho de la plantilla (7 y 8)
\newcommand{\apartado}[1]{%
  \par\vspace{14pt}%
  {\large\bfseries\color{urjc}#1\par}%
  \vspace{4pt}{\color{urjc}\hrule height 0.7pt}\vspace{5pt}}

% ---------- Recuadros ----------
% \begin{tarjeta}[subtítulo opcional]{título}
\newtcolorbox{tarjeta}[2][]{%
  enhanced, sharp corners, boxrule=0.5pt, titlerule=0pt,
  colframe=borde, colback=white, colbacktitle=urjc, coltitle=white,
  fonttitle=\bfseries\large, toptitle=3pt, bottomtitle=3pt,
  left=7pt, right=7pt, top=7pt, bottom=7pt,
  halign=flush left, halign title=flush left,
  before upper={\setlength{\parskip}{5pt}},
  title={#2\ifstrempty{#1}{}{\\[-1pt]{\footnotesize\itshape #1}}}}

% Fila de dos recuadros de igual altura
\newenvironment{filatarjetas}{%
  \begin{tcbraster}[raster columns=2, raster equal height=rows,
                    raster column skip=8pt, raster after skip=12pt]}
  {\end{tcbraster}}

% Sesión de la secuencia: \sesion{n}{tipo}{título}{duraciones y texto}
\newtcolorbox{sesionbox}[3]{%
  enhanced, sharp corners, boxrule=0.5pt, titlerule=0pt,
  colframe=borde, colback=white, colbacktitle=urjc, coltitle=white,
  fonttitle=\bfseries, toptitle=2pt, bottomtitle=2pt,
  left=6pt, right=6pt, top=4pt, bottom=5pt, after skip=8pt,
  halign=flush left, halign title=flush left,
  fontupper=\small, before upper={\setlength{\parskip}{3pt}},
  title={Sesión #1 · #2: #3}}

% Caja de información
\newtcolorbox{nota}[1]{%
  enhanced, sharp corners, boxrule=0pt, leftrule=3pt, colframe=urjc, colback=rosa,
  coltitle=urjc, colbacktitle=rosa, fonttitle=\bfseries, attach title to upper={\ },
  left=7pt, right=7pt, top=4pt, bottom=4pt, title={#1}, fontupper=\small}

% ---------- Tablas ----------
\NewTblrEnviron{tablaSA}
\SetTblrInner[tablaSA]{
  hlines={0.5pt,borde}, vlines={0.5pt,borde},
  leftsep=5pt, rightsep=5pt, rowsep=3pt,
  cells={font=\footnotesize, valign=t},
  row{1}={bg=rosa, fg=urjc, font=\footnotesize\bfseries},
}

% Versión con salto de página y cabecera repetida
\NewTblrTheme{limpio}{
  \SetTblrTemplate{caption}{empty}
  \SetTblrTemplate{capcont}{empty}
  \SetTblrTemplate{contfoot}{empty}
  \SetTblrTemplate{conthead}{empty}
  \SetTblrTemplate{note}{empty}
  \SetTblrTemplate{remark}{empty}
}
\NewTblrEnviron{tablaSAlarga}
\SetTblrOuter[tablaSAlarga]{long, theme=limpio, entry=none, label=none, presep=2pt, postsep=4pt}
\SetTblrInner[tablaSAlarga]{
  hlines={0.5pt,borde}, vlines={0.5pt,borde},
  leftsep=5pt, rightsep=5pt, rowsep=3pt,
  cells={font=\footnotesize, valign=t},
  rowhead=1,
}


\begin{document}


% ---------- Portada ----------
\thispagestyle{empty}
\vspace*{3.2cm}
\begin{center}
  {\small\color{pistagris}\MakeUppercase{\master}\\ Universidad Rey Juan Carlos}\\[1.4em]
  {\color{urjc}\rule{\linewidth}{1.2pt}}\\[1.2em]
  {\fontsize{13}{16}\selectfont\color{titular} Situación de aprendizaje}\\[0.8em]
  {\fontsize{25}{31}\selectfont\bfseries\color{titular}\tituloSA\par}
  \vspace{1em}
  {\color{urjc}\rule{\linewidth}{1.2pt}}\\[1.4em]
  {\large Matemáticas I · 1.º de Bachillerato (Ciencias y Tecnología)}\\[0.4em]
  {\large\color{pistagris}\asignatura}\\[3em]
  {\large\bfseries\autor}\\[0.4em]
  {\small\color{pistagris}Curso 2026-27}
\end{center}
\clearpage

{\setlength{\parskip}{0pt}\tableofcontents}
\clearpage

\section{Mapa de la situación de aprendizaje}

{\fontsize{22}{26}\selectfont\bfseries\color{titular}Situación de aprendizaje\par}
\vspace{6pt}

% ---------- 1 y 2 ----------
\begin{filatarjetas}
  \begin{tarjeta}{1 Identificación}
    \campo{Título:} \tituloSA

    \campo{Curso y materia:} 1.º de Bachillerato (Ciencias y Tecnología), Matemáticas I.

    \campo{N.º de sesiones / temporalización:} 20 sesiones de 55 minutos en 5 semanas (4 sesiones semanales). Es evaluable.

    \campo{Autor/a:} \autor

    \campo{Mapa (idea central):} con vectores, trigonometría y rectas podemos medir lo que no alcanzamos y decidir con datos dónde generar energía solar en el propio instituto.
  \end{tarjeta}
  \begin{tarjeta}{2 Contexto y reto}
    \campo{Situación / contexto:} el IES Ejemplo (centro ficticio de la Comunidad de Madrid) tiene paneles solares en la cubierta del gimnasio y quiere ampliar la instalación. Antes necesita saber qué sombras proyectan sus edificios, pero nadie ha medido las alturas y no se puede subir a la cubierta.

    \campo{Pregunta-reto:} ¿cómo podemos medir lo que no podemos alcanzar y usar esas medidas para decidir dónde caben más paneles sin que se hagan sombra?

    \campo{Lo que les cuento el primer día:} «La dirección os encarga un informe de verdad. Hace más de dos mil años, Eratóstenes midió la Tierra con una sombra; vosotros vais a medir el instituto con un clinómetro y a defender vuestra propuesta ante quien decide».
  \end{tarjeta}
\end{filatarjetas}

% ---------- 3 y 4 ----------
\begin{filatarjetas}
  \begin{tarjeta}[Matriz curricular 1]{3 Saberes básicos / criterios de evaluación}
    \begin{tablaSA}{colspec={X[1.5,l]X[1,l]}}
      Saberes básicos (bloque) & Criterios de evaluación \\
      Vectores: operaciones, módulo, ángulo, proyección (A) & 7.1, 2.1, 5.2 \\
      Razones, identidades y ecuaciones trigonométricas (B y D) & 3.2, 3.1, 1.2 \\
      Resolución de triángulos, teoremas del seno y del coseno (B) & 1.1, 1.2, 2.1 \\
      Ecuaciones de la recta, posición, ángulos y distancias (C) & 5.1, 5.2, 7.2 \\
      Modelización y herramientas digitales (C y D) & 1.1, 2.2, 7.2, 8.1, 6.1 \\
      Actitudes y aportación de las matemáticas (F) & 9.1, 9.2, 9.3, 6.2 \\
    \end{tablaSA}

    \pista{Texto literal y operativización completa en la sección~\ref{sec:matrices}.}
  \end{tarjeta}
  \begin{tarjeta}[Matriz curricular 2]{4 Competencias específicas y criterios}
    \begin{tablaSA}{colspec={Q[l,wd=30pt]X[l]}}
      CE & Criterios de evaluación \\
      CE1 & 1.1, 1.2: modelizar y obtener soluciones \\
      CE2 & 2.1, 2.2: verificar y elegir la solución \\
      CE3 & 3.1, 3.2: conjeturar con tecnología \\
      CE5 & 5.1, 5.2: conectar ideas matemáticas \\
      CE6 & 6.1, 6.2: conectar con otras áreas y retos \\
      CE7 & 7.1, 7.2: representar con tecnologías \\
      CE8 & 8.1: comunicar con rigor \\
      CE9 & 9.1, 9.2, 9.3: error, perseverancia, equipo \\
    \end{tablaSA}
  \end{tarjeta}
\end{filatarjetas}

% ---------- 5 y 6 ----------
\begin{filatarjetas}
  \begin{tarjeta}[Matriz curricular 3 · descriptores operativos]{5 Perfil de salida}
    \begin{tablaSA}{colspec={Q[l,wd=24pt]Q[l,wd=44pt]X[l]}}
      CE & Descriptor & Cómo se moviliza \\
      CE1 & STEM3 & Diseñar en equipo el modelo de sombras \\
      CE2 & STEM2 & Repetir medidas y valorar el error \\
      CE3 & STEM1 & De observar casos a demostrar \\
      CE5 & STEM1 & Resolver por varias vías conectadas \\
      CE6 & CC4 & Relacionar propuesta y sostenibilidad \\
      CE7 & CD2 & Crear el libro respetando la autoría \\
      CE8 & CCL1 & Defender la propuesta ante la dirección \\
      CE9 & CPSAA3.2 & Reparto equitativo de roles \\
    \end{tablaSA}
  \end{tarjeta}
  \begin{tarjeta}{6 Producto final}
    \campo{Producto:} libro de GeoGebra por equipo con la medición real, el modelo de sombras y una propuesta razonada de ampliación de paneles.

    \campo{Formato y destinatario:} libro digital y defensa oral de unos cinco minutos ante la clase y la dirección del centro, que recibe las propuestas. Los datos se comparten con el centro socio de eTwinning.

    \campo{Agrupamiento:} equipos heterogéneos de 4 estudiantes con roles rotatorios\soloProducto{ (sección~\ref{sec:cooperativo})}; trabajo individual en fichas, cuestionarios y examen.

    \campo{ODS:} 7 (energía asequible y no contaminante) y 11 (ciudades y comunidades sostenibles).
  \end{tarjeta}
\end{filatarjetas}

% ---------- 7 ----------
\apartado{7 Secuencia de actividades de aprendizaje}

\begin{tablaSAlarga}{
  colspec={Q[l,wd=44pt]X[2.2,l]X[1.1,l]X[1.3,l]X[1.4,l]},
  rowsep=3pt, cells={font=\scriptsize},
  column{1}={bg=rosa, fg=urjc, font=\scriptsize\bfseries},
  row{1}={bg=urjc, fg=white, font=\scriptsize\bfseries},
}
  Fase & Actividades / tareas & Sesiones y agrupamiento & Recursos & {Evaluación\\(instrumento y criterio)} \\
  Inicio & Presentación del reto y activación de conocimientos previos (triángulo rectángulo, Pitágoras y coordenadas). & S1. Gran grupo e individual & Plano del centro, GeoGebra Classroom & Cuestionario diagnóstico (sin calificación) \\
  Desarrollo & \textbf{Ejercicios:} tras cada exposición (vectores, trigonometría, triángulos, rectas), práctica guiada y una pieza del reto aplicada al propio centro. & S3, S4 y S6--S11. Magistral, individual y equipos & Pizarra, GeoGebra, hoja de ejercicios, plano del centro & Cuestionarios Q1--Q4 (7.1, 5.2, 1.2, 2.1) \\
  & \textbf{Actividades:} talleres de GeoGebra para explorar, conjeturar y aplicar: vectores, goniometría, y rayo y sombra. & S2, S5, S12 y S13 (F3, en S9). Individual y equipo & GeoGebra, fichas F1--F4, tablero Kanban & Fichas (7.1, 2.1, 3.1, 3.2, 1.1, 1.2, 7.2); mapa conceptual (5.1) \\
  & \textbf{Tareas:} medir el centro, construir el modelo de sombras y elaborar la propuesta de ampliación. & S14 a S17. Equipo & Clinómetro, cinta métrica, GeoGebra, PVGIS, eTwinning & Libro y propuesta (1.1, 2.2, 7.2, 6.1, 6.2); observación (9.1, 9.3) \\
  Cierre & Examen individual; defensa ante la dirección, coevaluación y metacognición. & S18 individual; S19 y S20 equipos & Examen modelo, libros, rúbricas & Examen; exposición (8.1, 1.1, 2.2, 6.1, 6.2); co/autoevaluación (7.2, 8.1, 9.x) \\
\end{tablaSAlarga}

% ---------- 8 ----------
\apartado{8 Evaluación}

\begin{tablaSAlarga}{
  colspec={X[1.2,l]X[2,l]X[1.3,l]}, cells={font=\scriptsize},
  row{1}={bg=rosa, fg=urjc, font=\scriptsize\bfseries},
}
  Criterios de evaluación & Instrumentos (rúbrica, lista de cotejo, diario…) & {Agente y momento\\(auto/co/hetero)} \\
  1.1, 1.2, 2.1, 3.1, 5.1, 5.2, 7.1, 7.2 & Examen escrito, cuestionarios y fichas GeoGebra & Heteroevaluación, durante los talleres y en S18 \\
  1.1, 2.2, 6.1, 6.2, 7.2, 8.1 & Rúbricas del libro y propuesta y de la exposición & Hetero y coevaluación, S16, S19 y S20 \\
  3.2, 5.1, 9.1, 9.2, 9.3 & Lista de cotejo, mapa conceptual y co/autoevaluación & Hetero, co y autoevaluación, S2--S20 \\
\end{tablaSAlarga}

\par\vspace{10pt}

% ---------- 9 ----------
\begin{tarjeta}{9 Atención a la diversidad (DUA)}
  Varias formas de acceder al contenido (explicación, construcción dinámica y material escrito), de expresarlo (libro, oral, escrito) y de implicarse (roles y elección del objeto que se mide). Refuerzo por criterio en S13 y ampliación con el radio de la Tierra. Detalle en la sección~\ref{sec:diversidad}.
\end{tarjeta}

% Versión con un centro y un grupo concretos: secciones/02_contexto_centro_concreto.tex
\section{Contexto y justificación}

\begin{nota}{Marco general.}
Esta situación de aprendizaje se integra en la programación general de Matemáticas~I. Por eso este apartado describe un marco común para un instituto público de la Comunidad de Madrid, y deja para cada centro la concreción de sus datos y de su grupo. El IES Ejemplo que aparece en el reto es un centro ficticio que sirve de ilustración.
\end{nota}

\subsection{Lugar en la programación}
La propuesta se sitúa en el segundo trimestre de Matemáticas~I, cuando se abordan la geometría y la trigonometría. Reúne en una sola unidad los vectores, la trigonometría y las rectas, y se planifica para 20 sesiones de 55 minutos, una duración que cada centro puede ajustar a su calendario.

\subsection{Condiciones del centro}
La propuesta puede aplicarse en cualquier instituto con Bachillerato de la modalidad de Ciencias y Tecnología. Solo requiere edificios u objetos altos que se puedan medir desde el exterior, un espacio al aire libre seguro, ordenadores con conexión a internet y permiso para trabajar sobre los espacios del propio centro. El reto gana sentido si el centro ya dispone de paneles solares o estudia instalarlos, pero también funciona como estudio de viabilidad. Cada centro adapta los objetos que se miden, los datos de latitud y las personas que reciben las propuestas.

\subsection{El alumnado}
La propuesta se dirige a estudiantes de 1.º de Bachillerato, habitualmente de 16 y 17 años. Se parte de un grupo heterogéneo en Matemáticas, donde convive quien domina los procedimientos pero tiene dificultad para justificarlos con quien necesita apoyo para pasar del triángulo rectángulo a los problemas generales. La programación prevé medidas generales para esa diversidad y orientaciones para las necesidades más habituales (sección~\ref{sec:diversidad}), que cada docente concreta según su grupo.

\subsection{Organización y tutoría}
Matemáticas I es una materia específica de modalidad con cuatro horas lectivas semanales \parencite[anexo~I, p.~32]{decreto64}. El departamento de Matemáticas coordina la propuesta dentro de su programación, y la comisión de coordinación pedagógica la articula con el resto de materias. El profesor tutor realiza el seguimiento del proceso educativo del grupo \parencite[art.~5.1, p.~17]{decreto64}, y el profesor de Matemáticas le informa del funcionamiento de los equipos. Como el horario del anexo~I no incluye una hora de tutoría en este curso \parencite[p.~32]{decreto64}, la formación en los roles se realiza dentro de las sesiones de Matemáticas\soloProducto{ (sección~\ref{sec:cooperativo})}. El departamento de orientación puede aportar información sobre salidas profesionales vinculadas a la energía, en línea con la finalidad orientadora del plan de acción tutorial \parencite[art.~5.3, p.~17]{decreto64}.

\subsection{Documentos del centro}
La propuesta debe ser coherente con los documentos institucionales de cada centro:
\begin{itemize}
  \item \textbf{Proyecto educativo:} recoge valores, fines y prioridades del centro e incorpora la concreción del currículo \parencite[art.~121.1, p.~122921]{lomloe}. La propuesta encaja especialmente si el centro prioriza la sostenibilidad.
  \item \textbf{Programación general anual:} el Consejo Escolar la aprueba y evalúa \parencite[art.~127.b, p.~122923]{lomloe}. Puede recoger la mejora energética del centro como objetivo del curso.
  \item \textbf{Plan de acción tutorial:} forma parte de la programación general anual \parencite[art.~5.3, p.~17]{decreto64}.
  \item \textbf{Programación didáctica de Matemáticas~I:} sitúa esta propuesta dentro de la secuencia anual de la materia.
\end{itemize}

\subsection{Marco legislativo}
La propuesta se rige por la Ley Orgánica 3/2020 \parencite{lomloe}, el Real Decreto 243/2022, que fija las enseñanzas mínimas y los descriptores del perfil de salida \parencite{rd243}, y el Decreto 64/2022 de la Comunidad de Madrid \parencite{decreto64}, cuyo anexo~II recoge las competencias específicas, los criterios y los contenidos de Matemáticas~I. El decreto pide planificar actividades que pongan en contexto la enseñanza y el aprendizaje \parencite[art.~16.3, p.~21]{decreto64} y permite a los centros adoptar innovaciones pedagógicas en el marco de su autonomía \parencite[art.~19.3.a, pp.~21--22]{decreto64}.

\subsection{Justificación de la propuesta}
\textbf{Por qué esta propuesta en este curso.} El currículo destaca el carácter instrumental de las matemáticas en otras áreas científicas y tecnológicas y sitúa la resolución de problemas y la investigación como aspectos fundamentales de la materia \parencite[p.~320]{decreto64}. Vectores, trigonometría y rectas suelen enseñarse por separado y este alumnado tiende a verlos sin relación. El problema de la sombra los necesita a todos: el rayo de sol es una recta con un vector director, su ángulo procede de la trigonometría y las alturas inaccesibles se obtienen resolviendo triángulos.

\textbf{Por qué con este enfoque.} Es frecuente que el alumnado de este curso domine los procedimientos, pero le cueste justificarlos. Defender una decisión real ante quien decide obliga a argumentar, y los equipos heterogéneos aprovechan la diversidad de niveles.

\subsection{Relación con otras materias}
\begin{itemize}
  \item \textbf{Tecnología e Ingeniería I:} energías renovables, eficiencia energética y sostenibilidad \parencite[p.~385]{decreto64}. Aporta los criterios técnicos de la instalación.
  \item \textbf{Física y Química:} carácter vectorial de las magnitudes y composición vectorial \parencite[p.~166]{decreto64}.
  \item \textbf{Orientación:} salidas profesionales vinculadas a la energía.
\end{itemize}

\subsection{Objetivos de Desarrollo Sostenible}
\begin{itemize}
  \item \textbf{ODS~7} \parencite{onu7}: la meta~7.2 propone aumentar la proporción de energía renovable. El alumnado estudia cómo aprovechar mejor una instalación solar real.
  \item \textbf{ODS~11} \parencite{onu11}: la meta~11.6 busca reducir el impacto ambiental de las ciudades. El instituto actúa como pequeña comunidad que planifica el uso de su propio espacio.
\end{itemize}

\section{Planificación y elementos curriculares}

\subsection{Objetivos didácticos}
Al finalizar, el alumnado será capaz de:
\begin{enumerate}[label=\textbf{O\arabic*.},leftmargin=2.4em]
  \item Representar vectores y operar con ellos en GeoGebra, interpretando geométricamente cada operación (CE7).
  \item Formular conjeturas sobre las razones trigonométricas y demostrarlas (CE3).
  \item Resolver triángulos con los teoremas del seno y del coseno y comprobar la validez de las soluciones (CE1, CE2).
  \item Relacionar vectores, trigonometría y geometría analítica en las ecuaciones de la recta, los ángulos y las distancias (CE5).
  \item Modelizar un problema real de medida y comunicar el método con rigor (CE1, CE7, CE8).
  \item Conectar el modelo con la energía solar y justificar una propuesta para el centro (CE2, CE6, CE9).
\end{enumerate}

\subsection{Competencias, criterios y saberes}
Se trabajan ocho de las nueve competencias específicas de Matemáticas~I (todas salvo la CE4), con sus criterios. De los descriptores del perfil de salida que el decreto conecta con cada competencia \parencite[pp.~321--324]{decreto64} se priorizan solo los que la secuencia trabaja con una actuación y una evidencia concretas. Los saberes son contenidos literales de los bloques A, B, C, D y F \parencite[pp.~325--328]{decreto64}. Las relaciones completas están en la sección~\ref{sec:matrices}.

\subsection{Metodología}
La secuencia combina tres modos de trabajo, cada uno con una función distinta:
\begin{itemize}
  \item \textbf{Sesiones magistrales (S3, S4 y S6--S11).} Cada bloque de contenidos se reparte en dos sesiones, con exposiciones de 20 a 30 minutos, para que demostrar y practicar quepa en el tiempo real de clase. El profesor expone de forma ordenada los contenidos, demuestra los resultados y resuelve ejemplos; después el alumnado practica con ejercicios y cierra con una \emph{pieza del reto}: diez minutos en equipo para aplicar lo aprendido al propio centro. Así el contexto está presente desde las primeras sesiones y no solo al final. El decreto recuerda que la ejercitación sigue siendo fundamental para aprender conceptos y métodos básicos antes de afrontar la resolución de problemas \parencite[p.~321]{decreto64}.
  \item \textbf{Talleres de GeoGebra (S2, S5 y S12).} El alumnado construye, manipula y conjetura. En vectores y trigonometría, el taller precede a la exposición, de modo que la sesión magistral formaliza lo observado. La herramienta actúa como colaboradora que ayuda a comprobar cálculos y conjeturas, no como caja negra que sustituye al cálculo \parencite[p.~321]{decreto64}.
  \item \textbf{Aprendizaje basado en un reto, con trabajo cooperativo (S1 y S14--S20).} Equipos con cuatro roles rotatorios (coordinación y Kanban, medición, modelo, comunicación)\soloProducto{, descritos en la sección~\ref{sec:cooperativo},} resuelven una tarea real. Esto responde al principio de fomentar la capacidad de aprender por sí mismos, de trabajar en equipo y de aplicar métodos de investigación apropiados \parencite[art.~4.1, p.~16]{decreto64}.
\end{itemize}
Los verbos de los criterios operativizados siguen la taxonomía revisada de Bloom \parencite{anderson2001} y avanzan desde la aplicación de procedimientos hasta la creación de un modelo propio.

\subsection{Temporalización}
Veinte sesiones de 55 minutos en cinco semanas, entre la segunda quincena de febrero y la tercera semana de marzo. La duración mínima de la guía es de tres semanas; aquí se amplía para que quepan los contenidos del bloque, el trabajo de campo y una defensa por equipo con tiempo suficiente.

\begin{tablaSAlarga}{colspec={Q[c,wd=26pt]Q[l,wd=52pt]Q[l,wd=84pt]X[l]Q[l,wd=44pt]},
  row{1}={bg=urjc, fg=white, font=\footnotesize\bfseries}}
  Ses. & Fase & Tipo & Contenido & Objetivos \\
  1 & Inicio & Reto & Reto, cuestionario diagnóstico, equipos y presentación del socio eTwinning & -- \\
  2 & Desarrollo & Taller (actividad) & Kanban; explorar vectores en GeoGebra & O1 \\
  3 & Desarrollo & \textbf{Magistral} + ejercicios & Vectores: estructura, módulo y coordenadas & O1 \\
  4 & Desarrollo & \textbf{Magistral} + ejercicios & Producto escalar, ángulo y proyección & O1 \\
  5 & Desarrollo & Taller (actividad) & Circunferencia goniométrica: conjeturas & O2 \\
  6 & Desarrollo & \textbf{Magistral} + ejercicios & Razones de un ángulo cualquiera, identidades, ángulo suma y diferencia & O2 \\
  7 & Desarrollo & \textbf{Magistral} + ejercicios & Ángulo doble y mitad; ecuaciones trigonométricas & O2 \\
  8 & Desarrollo & \textbf{Magistral} + ejercicios & Teoremas del seno y del coseno & O3 \\
  9 & Desarrollo & \textbf{Magistral} + ficha & Caso ambiguo y altura inaccesible & O3 \\
  10 & Desarrollo & \textbf{Magistral} + ejercicios & Ecuaciones de la recta y posiciones relativas & O4 \\
  11 & Desarrollo & \textbf{Magistral} + ejercicios & Ángulos entre rectas y distancias & O4 \\
  12 & Desarrollo & Taller (actividad) & Rayo y sombra en GeoGebra (F4) & O4, O5 \\
  13 & Desarrollo & Taller (actividad) & Refuerzo y ampliación por criterio; mapa conceptual; aprobación del método & O1--O5 \\
  14 & Desarrollo & Tarea & Medición en el patio al mediodía solar; datos eTwinning & O5, O6 \\
  15 & Desarrollo & Tarea & Modelo de sombras y análisis del error & O5, O6 \\
  16 & Desarrollo & Tarea & Propuesta de ampliación, comparación con el socio y coevaluación & O5, O6 \\
  17 & Desarrollo & Tarea & Mejora del libro y ensayo de la defensa & O5, O6 \\
  18 & Cierre & Evaluación & Examen individual y retroalimentación & O1--O4 \\
  19 & Cierre & Defensa & Defensa ante la dirección (cuatro equipos) y coevaluación & O5, O6 \\
  20 & Cierre & Defensa & Defensa (tres equipos), coevaluación, metacognición y cierre eTwinning & O5, O6 \\
\end{tablaSAlarga}

\section{Enfoque competencial: niveles de desempeño y evidencias}
La evaluación por competencias es el eje de la propuesta: cada actividad genera una evidencia y cada evidencia se valora con un nivel de desempeño. La tabla define cuatro niveles por competencia específica; las rúbricas de la sección~\ref{sec:instrumentos} los concretan.

\begin{tablaSAlarga}{colspec={Q[l,wd=30pt]X[l]X[l]X[l]X[l]Q[l,wd=74pt]},
  cells={font=\scriptsize}, colsep=3pt,
  row{1}={bg=urjc, fg=white, font=\scriptsize\bfseries},
  column{1}={font=\scriptsize\bfseries, fg=urjc}, cell{1}{1}={fg=white}}
  CE & 1 · Inicial & 2 · En desarrollo & 3 · Adecuado & 4 · Excelente & Evidencias \\
  CE1 & Aplica fórmulas sin relacionarlas con el problema & Modeliza con ayuda; omite alguna solución & Modeliza de forma autónoma y obtiene todas las soluciones & Además compara estrategias y justifica la más eficiente & Libro, F3, examen \\
  CE2 & No comprueba sus resultados & Comprueba con GeoGebra sin explicar diferencias & Comprueba y explica las discrepancias & Además estima el efecto del error y elige la solución según el contexto & F1, libro \\
  CE3 & Reproduce resultados dados & Formula conjeturas con guía, sin justificarlas & Formula conjeturas y las demuestra con ayuda & Demuestra de forma autónoma y propone casos nuevos & F2, examen \\
  CE5 & Trata cada contenido por separado & Relaciona dos ideas con ayuda & Conecta vectores, ángulos y rectas & Usa las conexiones para resolver de varias formas & Mapa, examen \\
  CE6 & No relaciona matemáticas y entorno & Señala alguna conexión general & Aplica el modelo a la instalación real & Además valora impacto y límites del modelo & Libro, defensa \\
  CE7 & Representaciones incorrectas o incompletas & Representa con una sola forma & Selecciona la representación adecuada & Justifica la elección y comunica con claridad & F1, F4, libro \\
  CE8 & Lenguaje impreciso y desordenado & Comunica con errores de notación & Comunica con orden y notación correcta & Además adapta el discurso a su público con rigor & Defensa \\
  CE9 & Evita tareas o abandona ante el error & Participa de forma irregular & Colabora y aprende de sus errores & Además impulsa al equipo e incorpora la crítica & Observación, co/auto \\
\end{tablaSAlarga}

\subsection{Correspondencia con la calificación}
Los niveles se traducen a la escala numérica así: \textbf{Inicial}, 1--4; \textbf{En desarrollo}, 5--6; \textbf{Adecuado}, 7--8; \textbf{Excelente}, 9--10. En las rúbricas R1 a R4 cada nivel vale un número fijo dentro de ese tramo: 2,5; 5; 7,5 y 10. La nota de cada instrumento se calcula como indica la sección~\ref{sec:calificacion}.

\subsection{Trabajo interdisciplinar}
La tarea no se puede resolver solo con Matemáticas: requiere criterios técnicos y físicos que aportan otras materias, y cada aportación deja una evidencia evaluable.

\begin{tablaSAlarga}{colspec={Q[l,wd=86pt]X[1.4,l]X[1.4,l]X[1,l]},
  row{1}={bg=urjc, fg=white, font=\footnotesize\bfseries}, column{1}={font=\footnotesize\bfseries}}
  Materia & Qué aporta & Dónde se trabaja & Evidencia \\
  Tecnología e Ingeniería I & Energías renovables y criterios de una instalación: orientación, separación entre filas y producción & S16, con la orientación del departamento de Tecnología sobre el pliego del IDAE & Plano y cálculo de filas en la propuesta (6.1) \\
  Física y Química & Magnitudes vectoriales y modelización del rayo como dirección & S4 y S12 & Ficha F4 y modelo de sombras (7.2) \\
  Centro y comunidad & Necesidad real, destinatario de la propuesta y datos del entorno, también del centro socio & S1, S14, S16, S19 y S20 & Defensa y reflexión sobre ODS (6.2) \\
\end{tablaSAlarga}

\clearpage
\begin{landscape}
\section{Matrices de relación curricular}\label{sec:matrices}
El criterio 7.1 no se evalúa en el examen, porque una prueba escrita no recoge la representación con tecnología; sus evidencias son la ficha F1 y el cuestionario Q1. Criterios y contenidos son literales del Decreto~64/2022 \parencite[pp.~324--328]{decreto64}.

Las relaciones se han depurado con un criterio: un descriptor solo se prioriza si hay una actuación concreta y una evidencia que respondan a su definición \parencite[anexo~I]{rd243}. El decreto conecta cada competencia con varios descriptores, pero esa conexión es general y no obliga a trabajarlos todos en una misma situación de aprendizaje \parencite[pp.~321--324]{decreto64}. Por eso las tablas indican también qué se deja fuera y por qué\soloJustificacion{, y explican cada relación entre criterio, sesión e instrumento}. Códigos de instrumentos: EX examen, CU cuestionario, FI fichas GeoGebra, LI libro y propuesta, EP exposición, MC mapa conceptual, OB observación, CO co/autoevaluación.
\subsection{Competencias específicas, criterios y descriptores}
\vspace{4pt}

\begin{tablaSAlarga}{colspec={Q[l,wd=8.2cm] Q[l,wd=1.5cm] Q[l,wd=3.7cm] Q[l,wd=9.6cm]},
  rowsep=2pt, colsep=4pt, cells={font=\scriptsize}, row{1}={bg=urjc, fg=white, font=\footnotesize\bfseries},}
Competencia específica & Criterios & Descriptores del decreto (en negrita, los priorizados) & Descriptores no priorizados y motivo \\
\textbf{CE1.} Modelizar y resolver problemas de la vida cotidiana y de la ciencia y la tecnología aplicando diferentes estrategias y formas de razonamiento para obtener posibles soluciones. & 1.1, 1.2 & STEM1, STEM2, \textbf{STEM3}, CD2, \textbf{CD5}, CPSAA4, CPSAA5, \textbf{CE3} & STEM1, STEM2 y CD2 se concretan en CE2, CE3 y CE7. CPSAA4 y CPSAA5 no tienen aquí una actividad específica. \\
\textbf{CE2.} Verificar la validez de las posibles soluciones de un problema empleando el razonamiento y la argumentación para contrastar su idoneidad. & 2.1, 2.2 & \textbf{STEM1}, \textbf{STEM2}, CD3, CPSAA4, CC3, CE3 & CPSAA4 se refiere a analizar información de medios de comunicación, y cotejar un cálculo no lo es. CD3, CC3 y CE3 no intervienen en la verificación. \\
\textbf{CE3.} Formular o investigar conjeturas o problemas, utilizando el razonamiento, la argumentación, la creatividad y el uso de herramientas tecnológicas, para generar nuevo conocimiento matemático. & 3.1, 3.2 & CCL1, \textbf{STEM1}, \textbf{STEM2}, CD1, CD2, CD3, CD5, CE3 & CD5 no se moviliza: al conjeturar, GeoGebra es un instrumento de investigación y no se crea ninguna herramienta. CD1, CD2, CD3, CE3 y CCL1 no definen lo que se aprende en esta tarea. \\
\textbf{CE5.} Establecer, investigar y utilizar conexiones entre las diferentes ideas matemáticas estableciendo vínculos entre conceptos, procedimientos, argumentos y modelos para dar significado y estructurar el aprendizaje matemático. & 5.1, 5.2 & \textbf{STEM1}, STEM3, CD2, CD3, CCEC1 & Compartir una construcción (CD3) o crear contenido digital (CD2) no conecta ideas matemáticas. STEM3 se concreta en CE1 y CCEC1 no tiene actividad. \\
\textbf{CE6.} Descubrir los vínculos de las matemáticas con otras áreas de conocimiento y profundizar en sus conexiones, interrelacionando conceptos y procedimientos, para modelizar, resolver problemas y desarrollar la capacidad crítica, creativa e innovadora en situaciones diversas. & 6.1, 6.2 & STEM1, \textbf{STEM2}, CD2, CPSAA5, \textbf{CC4}, CE2, CE3, CCEC1 & CE2 exigiría un análisis económico que la propuesta no incluye. CCEC1 se refiere al patrimonio cultural, y Eratóstenes solo se usa como motivación. STEM1, CD2 y CE3 se concretan en otras competencias; CPSAA5 no tiene actividad. \\
\textbf{CE7.} Representar conceptos, procedimientos e información matemáticos seleccionando diferentes tecnologías, para visualizar ideas y estructurar razonamientos matemáticos. & 7.1, 7.2 & STEM3, CD1, \textbf{CD2}, CD5, CE3, CCEC4.1, CCEC4.2 & CCEC4.1 y CCEC4.2 se refieren a proyectos artísticos, no al diseño de una construcción. CD5, STEM3 y CE3 se concretan una sola vez, en CE1. CD1 no tiene actividad específica. \\
\textbf{CE8.} Comunicar las ideas matemáticas, empleando el soporte, la terminología y el rigor apropiados, para organizar y consolidar el pensamiento matemático. & 8.1 & \textbf{CCL1}, CCL3, CP1, STEM2, \textbf{STEM4}, \textbf{CD3}, CCEC3.2 & CCL3 y CCEC3.2 no tienen actividad específica. CP1 no se evalúa, porque la lengua de trabajo con el socio no es objeto de esta materia. STEM2 se concreta en CE2 y CE3. \\
\textbf{CE9.} Utilizar destrezas personales y sociales, y organizando activamente el trabajo en equipos heterogéneos, aprendiendo del error como parte del proceso de aprendizaje y afrontando situaciones de incertidumbre, para perseverar en la consecución de objetivos en el aprendizaje de las matemáticas. & 9.1, 9.2, 9.3 & CP3, STEM5, \textbf{CPSAA1.1}, CPSAA1.2, CPSAA3.1, \textbf{CPSAA3.2}, CC2, CC3, CE2 & Los demás descriptores no tienen una actividad y una evidencia propias en esta secuencia. \\
\end{tablaSAlarga}

\subsection{Descriptores priorizados: actuación y justificación}
\vspace{4pt}

\begin{tablaSAlarga}{colspec={Q[l,wd=1cm] Q[l,wd=5.6cm] Q[l,wd=5.8cm] Q[l,wd=6.6cm] Q[l,wd=4cm]},
  rowsep=2pt, colsep=4pt, cells={font=\scriptsize}, row{1}={bg=urjc, fg=white, font=\footnotesize\bfseries},}
CE & Descriptor (fragmento literal) & Actuación concreta & Por qué la conexión es real & Evidencia \\
\textbf{CE1} & \textbf{STEM3.} «diseñando y creando prototipos o modelos para generar o utilizar productos que den solución a una necesidad o problema de forma colaborativa» & Cada equipo diseña el modelo de sombras y, con él, la propuesta de ampliación. & El reto es un proyecto: hay una necesidad real del centro y un producto que se evalúa según sus objetivos y su sostenibilidad. & Libro y propuesta (R1), S14--S16 \\
\textbf{CE1} & \textbf{CD5.} «Desarrolla soluciones tecnológicas innovadoras y sostenibles para dar respuesta a necesidades concretas» & El modelo de GeoGebra se entrega como herramienta que la dirección puede reutilizar cambiando la fecha y la altura del obstáculo. & Aquí sí se crea una herramienta digital que responde a una necesidad concreta. Por eso CD5 se trabaja en esta competencia y en ninguna otra. & Construcción dinámica del libro (R1, modelo y método), S15 \\
\textbf{CE1} & \textbf{CE3.} «estrategias ágiles de planificación y gestión de proyectos» & Cada equipo planifica el producto con un tablero Kanban y lo revisa en cada sesión de trabajo. & Kanban es una estrategia ágil, y S14--S17 obligan a gestionar tareas y tiempo para llegar a un resultado útil para otros. & Tablero Kanban (L1), S2--S17 \\
\textbf{CE2} & \textbf{STEM1.} «analizando críticamente las soluciones y reformulando el procedimiento» & El alumnado compara el cálculo a mano con GeoGebra, explica la discrepancia y rehace el procedimiento si hace falta. & El descriptor describe lo mismo que pide la competencia: analizar una solución y corregir el camino seguido. & F1, F3 y examen, pregunta 2 \\
\textbf{CE2} & \textbf{STEM2.} «apreciando la importancia de la precisión y la veracidad» & Cada medida se repite y se calcula cómo cambia la altura con un grado de error. & A diferencia de un problema de libro, la medida real tiene error, y de él depende la validez de la solución. & Análisis del error en el libro (R1, error y solución), S14--S15 \\
\textbf{CE3} & \textbf{STEM1.} «métodos inductivos y deductivos propios del razonamiento matemático» & En F2 se observan casos y se formula una conjetura (S5); después se demuestra (S6). & La secuencia reproduce los dos métodos que nombra el descriptor: primero inducción y luego deducción. & F2 y examen, pregunta 3 \\
\textbf{CE3} & \textbf{STEM2.} «planteándose hipótesis y contrastándolas o comprobándolas mediante la observación» & En F2 se pone a prueba en GeoGebra la conjetura $\sin 2\alpha = 2\sin\alpha$ y se busca un contraejemplo. & La conjetura funciona como hipótesis y GeoGebra como instrumento de observación, no como producto. & F2 y observación (L1), S5 \\
\textbf{CE5} & \textbf{STEM1.} «emplea estrategias variadas para la resolución de problemas» & Ángulos, perpendicularidad y distancias se abordan con producto escalar, ecuaciones de la recta y trigonometría. & Conectar ideas se demuestra al poder resolver un problema por más de una vía y elegir entre ellas. & Mapa conceptual (L2) y examen, pregunta 5 \\
\textbf{CE6} & \textbf{CC4.} «actividades y hábitos que conduzcan al logro de los Objetivos de Desarrollo Sostenible» & Cada equipo argumenta cómo contribuye su propuesta a los ODS 7 y 11. & La decisión afecta a la energía de su propio centro: una actuación local conectada con un objetivo global. & Reflexión sobre los ODS en el libro (R1, ODS y límites) y defensa (R2) \\
\textbf{CE6} & \textbf{STEM2.} «mostrando una actitud crítica acerca del alcance y limitaciones de los métodos empleados» & Se compara la separación del mediodía con la del pliego del IDAE y se señalan los límites del modelo. & El modelo matemático por sí solo no basta: el criterio técnico muestra hasta dónde llega y dónde se queda corto. & Propuesta y límites del modelo en el libro (R1, propuesta y ODS y límites) \\
\textbf{CE7} & \textbf{CD2.} «Crea, integra y reelabora contenidos digitales de forma individual o colectiva» & El equipo integra en el libro construcciones, cálculos y explicación, y cita los datos y materiales ajenos. & El libro es un contenido digital colectivo que reelabora datos de terceros, por lo que también obliga a respetar su autoría. & Construcciones y fuentes citadas en el libro (R1, representación) \\
\textbf{CE8} & \textbf{CCL1.} «adecuación a los diferentes contextos sociales y académicos» & La propuesta se defiende ante la clase y la dirección. & Explicar el mismo método a compañeros y a quien decide obliga a adaptar el discurso sin perder rigor. & Exposición (R2), S19--S20 \\
\textbf{CE8} & \textbf{STEM4.} «de forma clara y precisa, en diferentes formatos» & El método y los resultados se presentan con fórmulas, tablas, gráficos y construcciones. & Cada formato aporta algo distinto: el plano convence a la dirección y la fórmula justifica el cálculo. & Exposición (R2) y libro (R1) \\
\textbf{CE8} & \textbf{CD3.} «para comunicarse, trabajar colaborativamente y compartir información, gestionando de manera responsable sus acciones, presencia y visibilidad en la red» & Datos y conclusiones se comparten en el TwinSpace con avatares y sin datos personales. & Es un intercambio real con otro centro, en el que la identidad digital se gestiona con reglas explícitas. & Publicaciones del TwinSpace (seguimiento, sin calificación) \\
\textbf{CE9} & \textbf{CPSAA3.2.} «Distribuye en un grupo las tareas, recursos y responsabilidades de manera ecuánime» & Los cuatro roles rotan, de modo que todos los ejercen todos\soloProducto{ (sección~\ref{sec:cooperativo})}. & La rotación convierte el reparto ecuánime en una regla de funcionamiento, no en una intención. & Tablero Kanban y cumplimiento del rol (L1), R3 \\
\textbf{CE9} & \textbf{CPSAA1.1.} «la resiliencia, la autoeficacia y la búsqueda de objetivos de forma autónoma» & Las conjeturas falsas y las medidas fallidas se revisan y se repiten. & El reto incluye fallos reales (nubes, medidas incoherentes), que dan ocasión de perseverar. & Observación (L1) y autoevaluación (R3) \\
\end{tablaSAlarga}

\subsection{Criterios operativizados, saberes, sesiones, instrumentos y evidencias}
\vspace{4pt}

\begin{tablaSAlarga}{colspec={Q[l,wd=5.2cm] Q[l,wd=7.2cm] Q[c,wd=1.3cm] Q[c,wd=1.3cm] Q[c,wd=1.5cm] Q[l,wd=6.5cm]},
  rowsep=2pt, colsep=4pt, cells={font=\scriptsize}, row{1}={bg=urjc, fg=white, font=\footnotesize\bfseries},}
Criterio de evaluación & Criterio operativizado & Bloom & Sesiones & Instrum. & Evidencia de aprendizaje \\
\SetCell[c=6]{l,bg=rosa} {\textbf{A. Números y operaciones: vectores}\\ \textbullet~Adición y producto escalar de vectores: propiedades y representaciones.\\ \textbullet~Conjunto de vectores: estructura, comprensión y propiedades.\\ \textbullet~Módulo de un vector, coordenadas de un vector con respecto a una base, ángulo entre dos vectores y proyección ortogonal.\\ \textbullet~Estrategias para operar con números reales y vectores: cálculo mental o escrito en los casos sencillos y con herramientas tecnológicas en los casos más complicados.} & & & & & \\*
\textbf{7.1.} Representar ideas matemáticas, estructurando diferentes razonamientos matemáticos y seleccionando las tecnologías más adecuadas. & El alumno ha representado vectores en GeoGebra y ha calculado su suma, su producto por un número y su producto escalar, interpretando geométricamente cada resultado. & \textit{Aplicar} & 2, 3, 4 & CU, FI & Ficha F1 y cuestionario Q1 con operaciones interpretadas. \\
\textbf{2.1.} Comprobar la validez matemática de las posibles soluciones de un problema utilizando el razonamiento y la argumentación. & El alumno ha comprobado con GeoGebra el módulo, el ángulo entre dos vectores y la proyección ortogonal calculados a mano, justificando el origen de las discrepancias. & \textit{Evaluar} & 2, 4, 18 & EX, FI & Ficha F1 con cálculo manual, comprobación y causa de la discrepancia. \\
\textbf{5.2.} Resolver problemas en contextos matemáticos estableciendo y aplicando conexiones entre las diferentes ideas matemáticas. & El alumno ha resuelto problemas de perpendicularidad y de cálculo de ángulos aplicando el producto escalar y relacionando su expresión en coordenadas con su significado geométrico. & \textit{Aplicar} & 4, 18 & EX, CU & Problemas de perpendicularidad y ángulos resueltos. \\
\SetCell[c=6]{l,bg=rosa} {\textbf{B y D. Trigonometría}\\ \textbullet~Razones trigonométricas de un ángulo cualquiera medido en grados o en radianes.\\ \textbullet~Demostración de las identidades trigonométricas. Razones trigonométricas del ángulo suma, el ángulo diferencia, el ángulo doble y el ángulo mitad.\\ \textbullet~Cálculo de las razones trigonométricas de un ángulo cualquiera empleando las principales fórmulas trigonométricas.\\ \textbullet~Conjeturas geométricas en el plano: validación por medio de la deducción y la demostración de teoremas.\\ \textbullet~Resolución de ecuaciones trigonométricas sencillas.} & & & & & \\*
\textbf{3.2.} Emplear herramientas tecnológicas adecuadas en la formulación o investigación de conjeturas o problemas. & El alumno ha examinado en una circunferencia goniométrica dinámica de GeoGebra (con un deslizador para el ángulo) cómo varían el seno, el coseno y la tangente en los cuatro cuadrantes. & \textit{Analizar} & 5 & FI, OB & Ficha F2 con la tabla de signos y observaciones por cuadrante. \\
\textbf{3.1.} Adquirir nuevo conocimiento matemático a partir de la formulación de conjeturas y problemas de forma guiada. & El alumno ha formulado, de forma guiada, conjeturas sobre identidades trigonométricas (fundamental, ángulo doble, ángulo suma) a partir de la construcción y las ha demostrado por escrito. & \textit{Crear} & 5, 6, 18 & EX, FI & Conjetura escrita en F2 y demostración en papel. \\
\textbf{1.2.} Obtener todas las posibles soluciones matemáticas de problemas de la vida cotidiana y de la ciencia y la tecnología, describiendo el procedimiento utilizado. & El alumno ha resuelto ecuaciones trigonométricas sencillas obteniendo todas sus soluciones, tanto en [0°, 360°) como en forma general, y ha descrito el procedimiento seguido. & \textit{Aplicar} & 7, 18 & EX, CU & Cuestionario Q2: soluciones en el intervalo y solución general. \\
\SetCell[c=6]{l,bg=rosa} {\textbf{B. Medida y geometría: resolución de triángulos}\\ \textbullet~Cálculo de longitudes y medidas angulares: uso de la trigonometría.\\ \textbullet~Aplicación de las razones trigonométricas, el teorema de los senos y el teorema del coseno en la resolución de triángulos y de problemas geométricos de contexto real.\\ \textbullet~Demostración del teorema del seno y del coseno.} & & & & & \\*
\textbf{1.1.} Manejar algunas estrategias y herramientas, incluidas las digitales, en la modelización y resolución de problemas de la vida cotidiana y de la ciencia y la tecnología, evaluando su eficiencia en cada caso. & El alumno ha justificado la elección entre el teorema del seno y el del coseno según los datos de cada triángulo, valorando la eficiencia de cada estrategia. & \textit{Evaluar} & 8, 9, 18 & EX, FI & Ficha F3: elección razonada del teorema. \\
\textbf{1.2.} Obtener todas las posibles soluciones matemáticas de problemas de la vida cotidiana y de la ciencia y la tecnología, describiendo el procedimiento utilizado. & El alumno ha discriminado si un triángulo dado por dos lados y el ángulo opuesto a uno de ellos tiene dos, una o ninguna solución y ha calculado todas las soluciones, describiendo el procedimiento. & \textit{Analizar} & 9, 18 & EX, FI & Ficha F3: número de soluciones y procedimiento. \\
\textbf{2.1.} Comprobar la validez matemática de las posibles soluciones de un problema utilizando el razonamiento y la argumentación. & El alumno ha verificado en GeoGebra, arrastrando los vértices, la validez de los triángulos obtenidos analíticamente y ha argumentado por qué un caso no tiene solución. & \textit{Evaluar} & 9 & CU, FI & Cuestionario Q3 y comprobación arrastrando vértices. \\
\SetCell[c=6]{l,bg=rosa} {\textbf{C. Geometría en el plano: rectas}\\ \textbullet~Ecuaciones de la recta en el espacio bidimensional.\\ \textbullet~Estudio de la posición relativa de puntos y rectas en el plano.\\ \textbullet~Planteamiento y resolución de problemas de geometría afín relacionados con la incidencia, el paralelismo y la ortogonalidad de rectas en el plano.\\ \textbullet~Planteamiento y resolución de problemas de geometría métrica relacionados con la medida de ángulos entre rectas y la medida de distancias entre puntos y rectas.\\ \textbullet~Expresiones algebraicas de objetos geométricos: selección de la más adecuada en función de la situación a resolver.} & & & & & \\*
\textbf{5.1.} Manifestar una visión matemática integrada, investigando y conectando las diferentes ideas matemáticas. & El alumno ha relacionado el vector director y el vector normal con las distintas ecuaciones de la recta, organizando en un mapa conceptual las conexiones entre vectores, trigonometría y geometría analítica. & \textit{Analizar} & 10, 13, 18 & EX, MC & Mapa conceptual con enlaces entre vector director, normal y ángulo. \\
\textbf{5.2.} Resolver problemas en contextos matemáticos estableciendo y aplicando conexiones entre las diferentes ideas matemáticas. & El alumno ha determinado la posición relativa de dos rectas y ha calculado ángulos entre rectas y distancias de un punto a una recta aplicando el producto escalar. & \textit{Aplicar} & 11, 18 & EX, CU & Cuestionario Q4 con posición, ángulo y distancia. \\
\textbf{7.2.} Seleccionar y utilizar diversas formas de representación, valorando su utilidad para compartir información. & El alumno ha seleccionado la ecuación de la recta (vectorial, paramétrica, continua, general o explícita) más útil para cada problema, justificando su elección. & \textit{Evaluar} & 10, 12, 18 & EX, FI & Ficha F4: ecuación elegida y motivo. \\
\SetCell[c=6]{l,bg=rosa} {\textbf{C y D. Modelización: producto final «Medir lo que no se puede alcanzar»}\\ \textbullet~Representación de objetos geométricos en el plano mediante herramientas digitales.\\ \textbullet~Modelos matemáticos (geométricos, algebraicos, grafos...) en la resolución de problemas en el plano. Conexiones con otras disciplinas y áreas de interés.\\ \textbullet~Modelización de la posición y el movimiento de un objeto en el plano mediante vectores.\\ \textbullet~Formulación, resolución y análisis de problemas de la vida cotidiana y de la ciencia y la tecnología utilizando herramientas o programas adecuados.} & & & & & \\*
\textbf{1.1.} Manejar algunas estrategias y herramientas, incluidas las digitales, en la modelización y resolución de problemas de la vida cotidiana y de la ciencia y la tecnología, evaluando su eficiencia en cada caso. & El alumno ha diseñado un método de medida y ha construido un modelo a escala en GeoGebra para calcular una altura o distancia inaccesible del centro, valorando su eficiencia frente a otros métodos. & \textit{Crear} & 14, 15, 19, 20 & LI, EP & Libro: método, construcción y comparación con otro método. \\
\textbf{2.2.} Seleccionar la solución más adecuada de un problema en función del contexto, usando el razonamiento y la argumentación. & El alumno ha valorado el efecto de un error de un grado en el resultado final y ha seleccionado, entre las medidas repetidas, el resultado más razonable. & \textit{Evaluar} & 14, 15, 19, 20 & LI, EP & Libro: medidas repetidas, sensibilidad y resultado elegido. \\
\textbf{7.2.} Seleccionar y utilizar diversas formas de representación, valorando su utilidad para compartir información. & El alumno ha elaborado un libro de GeoGebra que integra la construcción, los cálculos y la explicación del método, eligiendo las representaciones más claras para compartirlo con la clase. & \textit{Crear} & 15, 16, 17, 19, 20 & LI, CO & Libro completo con representaciones justificadas. \\
\textbf{8.1.} Mostrar organización al comunicar las ideas matemáticas empleando el soporte, la terminología y el rigor apropiados. & El alumno ha explicado oralmente el método y los resultados con orden, utilizando la terminología y la notación matemáticas con rigor. & \textit{Comprender} & 17, 19, 20 & EP, CO & Ensayo y defensa oral con terminología y notación correctas. \\
\SetCell[c=6]{l,bg=rosa} {\textbf{F. Actitudes y aprendizaje}\\ \textbullet~Tratamiento del error como elemento movilizador de conocimientos previos adquiridos y generador de oportunidades de aprendizaje en el aula de matemáticas.\\ \textbullet~Reconocimiento y aceptación de diversos planteamientos en la resolución de problemas, transformando los enfoques de los demás en nuevas y mejoradas estrategias propias.\\ \textbullet~Técnicas y estrategias de trabajo en equipo para la resolución de problemas y tareas matemáticas, en grupos heterogéneos.\\ \textbullet~Destrezas para desarrollar una comunicación efectiva: la escucha activa, la formulación de preguntas o solicitud y prestación de ayuda cuando sea necesario.} & & & & & \\*
\textbf{9.1.} Afrontar las situaciones de incertidumbre, aceptando y aprendiendo del error como parte del proceso de aprendizaje de las matemáticas. & El alumno ha analizado los errores de sus conjeturas y de sus medidas, y los ha utilizado como punto de partida para revisar su razonamiento en la puesta en común. & \textit{Analizar} & 5, 12, 14, 19, 20 & OB, CO & Registro de observación y revisión de medidas fallidas. \\
\textbf{9.2.} Mostrar una actitud positiva y perseverante, aceptando y aprendiendo de la crítica razonada al hacer frente a las diferentes actividades de las matemáticas. & El alumno ha valorado la crítica razonada recibida en la coevaluación y la ha incorporado para mejorar el libro de GeoGebra, perseverando ante las dificultades. & \textit{Evaluar} & 16, 17, 19, 20 & OB, CO & Registro de cambios del libro tras la coevaluación. \\
\textbf{9.3.} Participar en tareas matemáticas de forma activa en equipos heterogéneos, escuchando su razonamiento. & El alumno ha colaborado activamente en el equipo heterogéneo, ha asumido tareas del tablero Kanban y ha tenido en cuenta los razonamientos de sus compañeros. & \textit{Aplicar} & 2, 14, 15, 16, 19, 20 & OB, CO & Tablero Kanban, roles rotados y registro de observación. \\
\SetCell[c=6]{l,bg=rosa} {\textbf{C y F. Conexión con la energía solar}\\ \textbullet~Modelos matemáticos (geométricos, algebraicos, grafos...) en la resolución de problemas en el plano. Conexiones con otras disciplinas y áreas de interés.\\ \textbullet~Valoración de la contribución de las matemáticas y el papel de matemáticos a lo largo de la historia en el avance de la ciencia y la tecnología.} & & & & & \\*
\textbf{6.1.} Resolver problemas en situaciones diversas, utilizando procesos matemáticos, estableciendo y aplicando conexiones entre el mundo real, otras áreas de conocimiento y las matemáticas. & El alumno ha resuelto la ubicación de los nuevos paneles conectando las medidas reales, el modelo geométrico y el criterio técnico de separación entre filas. & \textit{Aplicar} & 14, 16, 19, 20 & LI, EP & Propuesta de ampliación con plano, cálculo de filas y decisión justificada. \\
\textbf{6.2.} Analizar la aportación de las matemáticas al progreso de la humanidad reflexionando sobre su contribución en la propuesta de soluciones a situaciones complejas y a los retos científicos y tecnológicos que se plantean en la sociedad. & El alumno ha analizado cómo contribuye su propuesta a los retos de energía sostenible (ODS 7 y 11), señalando el papel del modelo matemático y sus límites. & \textit{Analizar} & 16, 19, 20 & LI, EP & Apartado del libro sobre ODS, límites del modelo y defensa. \\
\end{tablaSAlarga}

\ifjustificacion
\subsection{Justificación pedagógica de las relaciones}
\vspace{4pt}

\begin{tablaSAlarga}{colspec={Q[l,wd=1cm] Q[l,wd=3.4cm] Q[c,wd=1.4cm] Q[l,wd=8.9cm] Q[c,wd=1.3cm] Q[l,wd=7cm]},
  rowsep=2pt, colsep=4pt, cells={font=\scriptsize}, row{1}={bg=urjc, fg=white, font=\footnotesize\bfseries},}
Crit. & Aspecto & Sesiones & Por qué en esas sesiones & Instrum. & Por qué ese instrumento \\
\textbf{7.1} & Representar y operar con vectores & 2, 3, 4 & Se construye primero en F1 (S2), para que cada operación tenga un significado geométrico antes de formalizarla en S3 y S4. & CU, FI & La ficha recoge la representación con tecnología y Q1 comprueba su interpretación. El examen escrito no permite usar GeoGebra. \\
\textbf{2.1} & Validar cálculos vectoriales & 2, 4, 18 & Validar exige tener un cálculo propio: en S2 se contrasta de forma exploratoria y en S4, con la fórmula de la proyección, se explica la discrepancia. S18 comprueba si se hace sin herramienta. & EX, FI & F1 registra el contraste entre mano y GeoGebra; la pregunta 2 del examen obliga a detectar un error ajeno. \\
\textbf{5.2} & Producto escalar en problemas & 4, 18 & El producto escalar une coordenadas y geometría, y se enseña en S4; S18 comprueba que se aplica de forma autónoma. & EX, CU & Q1 da un diagnóstico temprano para el refuerzo por criterio de S13, y el examen verifica la transferencia. \\
\textbf{3.2} & Investigar con la circunferencia & 5 & S5 precede a propósito a la exposición de S6: hay que investigar antes de conocer la fórmula. & FI, OB & La ficha muestra lo investigado; la observación capta si se explora antes de pedir ayuda, que no queda en el papel. \\
\textbf{3.1} & Conjeturar y demostrar & 5, 6, 18 & La conjetura nace en S5 y se demuestra en S6, de modo que la demostración responde a una pregunta propia. S18 comprueba que se demuestra sin apoyo. & EX, FI & F2 recoge la conjetura y la pregunta 3a del examen, la demostración individual. \\
\textbf{1.2} & Ecuaciones trigonométricas & 7, 18 & Las ecuaciones usan el ángulo doble, que se estudia en la misma sesión, y la pieza del reto les da un uso inmediato. & EX, CU & Obtener todas las soluciones es un procedimiento: Q2 da retorno inmediato y el examen lo certifica. \\
\textbf{1.1} & Elegir el teorema adecuado & 8, 9, 18 & La elección se aprende en S8 con datos variados y se pone a prueba en S9, con el caso que genera más dudas. & EX, FI & F3 y la pregunta 4 piden justificar la elección, no solo dar el resultado. \\
\textbf{1.2} & Caso ambiguo & 9, 18 & Va al final del bloque porque exige dominar antes el teorema del seno. & EX, FI & En F3 se ve cuántas soluciones hay al arrastrar; la pregunta 4b exige obtenerlas por cálculo. \\
\textbf{2.1} & Validar triángulos & 9 & S9 concentra el caso en que validar es decisivo: si $\sin B>1$, el triángulo no existe. & CU, FI & La comprobación arrastrando vértices solo se ve en F3; Q3 la cierra con un caso sin solución. \\
\textbf{5.1} & Conectar vector y recta & 10, 13, 18 & S10 deduce las ecuaciones del vector director; el mapa llega en S13, cuando ya se han visto los tres bloques y hay algo que conectar. & EX, MC & El mapa hace visibles las conexiones y la pregunta 5 del examen comprueba que se usan. \\
\textbf{5.2} & Ángulos y distancias & 11, 18 & Reutilizan el producto escalar de S4; por eso se sitúan al final del bloque, para volver sobre él. & EX, CU & Q4 alimenta el refuerzo por criterio de S13 y el examen verifica el dominio individual. \\
\textbf{7.2} & Elegir la ecuación de la recta & 10, 12, 18 & S10 presenta las cinco formas y F4, en S12, obliga a elegir una para un problema real. & EX, FI & F4 y la pregunta 5d piden explicar por qué se elige esa ecuación. \\
\textbf{1.1} & Método y modelo & 14, 15, 19, 20 & Se aplica cuando ya se dispone de todas las herramientas: el método se aprueba en S13, se ejecuta en S14 y se modela en S15. & LI, EP & El libro documenta el proceso; la defensa comprueba que todo el equipo lo entiende. \\
\textbf{2.2} & Elegir el resultado & 14, 15, 19, 20 & Elegir solo tiene sentido con medidas reales y repetidas (S14), y la sensibilidad al error se analiza con ellas en S15. & LI, EP & El libro recoge datos y decisión; en la defensa hay que justificarla ante preguntas. \\
\textbf{7.2} & Integrar representaciones & 15, 16, 17, 19, 20 & Integrar exige tener ya datos y modelo (S15 y S16), y S17 permite mejorarlo; las defensas muestran si otras personas lo entienden. & LI, CO & El libro es la representación. La coevaluación aporta lo que pide el criterio: su utilidad para compartir información. \\
\textbf{8.1} & Comunicar con rigor & 17, 19, 20 & Se ensaya en S17 y se evalúa en S19 y S20, ante un público real. & EP, CO & La comunicación oral solo se puede evaluar en una exposición; la co/autoevaluación añade la mirada del público. \\
\textbf{9.1} & Aprender del error & 5, 12, 14, 19, 20 & Se observa donde el error aparece de forma natural: conjeturas falsas (S5), construcciones (S12) y medidas (S14). & OB, CO & Una actitud no cabe en un examen: se observa cuando ocurre y se reflexiona en la autoevaluación. \\
\textbf{9.2} & Aceptar la crítica & 16, 17, 19, 20 & Necesita una crítica aprovechable: la coevaluación del borrador en S16 se aprovecha en S17, dentro del horario de clase. & OB, CO & El registro de cambios permite ver qué se ha incorporado realmente. \\
\textbf{9.3} & Colaborar en equipo & 2, 14, 15, 16, 19, 20 & Los roles se presentan en S2 y se exigen sobre todo de S14 a S16, cuando el producto depende de todos. & OB, CO & Solo la observación y los compañeros pueden valorar la colaboración efectiva. \\
\textbf{6.1} & Conectar con la instalación & 14, 16, 19, 20 & La conexión con el criterio técnico requiere haber medido; se hace en S16, con la orientación del departamento de Tecnología. & LI, EP & La propuesta recoge el cálculo y la defensa, su viabilidad ante quien decide. \\
\textbf{6.2} & Valorar impacto y límites & 16, 19, 20 & La reflexión tiene sentido cuando el modelo ya existe y se conocen sus resultados. & LI, EP & El apartado del libro la deja por escrito y la defensa la somete a preguntas. \\
\end{tablaSAlarga}

\fi

\clearpage
\subsection{Cobertura: cada criterio por sesiones e instrumentos}
\vspace{4pt}

\begin{tablaSAlarga}{colspec={Q[c,wd=1.3cm] *{20}{Q[c,wd=0.58cm]} *{8}{Q[c,wd=0.75cm]}},
  rowsep=1.5pt, colsep=2pt, cells={font=\scriptsize}, row{1}={bg=urjc, fg=white, font=\scriptsize\bfseries},
  column{22-29}={bg=rosa}, cell{1}{22-29}={bg=urjc},}
Criterio & S1 & S2 & S3 & S4 & S5 & S6 & S7 & S8 & S9 & S10 & S11 & S12 & S13 & S14 & S15 & S16 & S17 & S18 & S19 & S20 & EX & CU & FI & LI & EP & MC & OB & CO \\
\textbf{1.1} &  &  &  &  &  &  &  & \textbf{X} & \textbf{X} &  &  &  &  & \textbf{X} & \textbf{X} &  &  & \textbf{X} & \textbf{X} & \textbf{X} & \textbf{X} &  & \textbf{X} & \textbf{X} & \textbf{X} &  &  &  \\
\textbf{1.2} &  &  &  &  &  &  & \textbf{X} &  & \textbf{X} &  &  &  &  &  &  &  &  & \textbf{X} &  &  & \textbf{X} & \textbf{X} & \textbf{X} &  &  &  &  &  \\
\textbf{2.1} &  & \textbf{X} &  & \textbf{X} &  &  &  &  & \textbf{X} &  &  &  &  &  &  &  &  & \textbf{X} &  &  & \textbf{X} & \textbf{X} & \textbf{X} &  &  &  &  &  \\
\textbf{2.2} &  &  &  &  &  &  &  &  &  &  &  &  &  & \textbf{X} & \textbf{X} &  &  &  & \textbf{X} & \textbf{X} &  &  &  & \textbf{X} & \textbf{X} &  &  &  \\
\textbf{3.1} &  &  &  &  & \textbf{X} & \textbf{X} &  &  &  &  &  &  &  &  &  &  &  & \textbf{X} &  &  & \textbf{X} &  & \textbf{X} &  &  &  &  &  \\
\textbf{3.2} &  &  &  &  & \textbf{X} &  &  &  &  &  &  &  &  &  &  &  &  &  &  &  &  &  & \textbf{X} &  &  &  & \textbf{X} &  \\
\textbf{5.1} &  &  &  &  &  &  &  &  &  & \textbf{X} &  &  & \textbf{X} &  &  &  &  & \textbf{X} &  &  & \textbf{X} &  &  &  &  & \textbf{X} &  &  \\
\textbf{5.2} &  &  &  & \textbf{X} &  &  &  &  &  &  & \textbf{X} &  &  &  &  &  &  & \textbf{X} &  &  & \textbf{X} & \textbf{X} &  &  &  &  &  &  \\
\textbf{6.1} &  &  &  &  &  &  &  &  &  &  &  &  &  & \textbf{X} &  & \textbf{X} &  &  & \textbf{X} & \textbf{X} &  &  &  & \textbf{X} & \textbf{X} &  &  &  \\
\textbf{6.2} &  &  &  &  &  &  &  &  &  &  &  &  &  &  &  & \textbf{X} &  &  & \textbf{X} & \textbf{X} &  &  &  & \textbf{X} & \textbf{X} &  &  &  \\
\textbf{7.1} &  & \textbf{X} & \textbf{X} & \textbf{X} &  &  &  &  &  &  &  &  &  &  &  &  &  &  &  &  &  & \textbf{X} & \textbf{X} &  &  &  &  &  \\
\textbf{7.2} &  &  &  &  &  &  &  &  &  & \textbf{X} &  & \textbf{X} &  &  & \textbf{X} & \textbf{X} & \textbf{X} & \textbf{X} & \textbf{X} & \textbf{X} & \textbf{X} &  & \textbf{X} & \textbf{X} &  &  &  & \textbf{X} \\
\textbf{8.1} &  &  &  &  &  &  &  &  &  &  &  &  &  &  &  &  & \textbf{X} &  & \textbf{X} & \textbf{X} &  &  &  &  & \textbf{X} &  &  & \textbf{X} \\
\textbf{9.1} &  &  &  &  & \textbf{X} &  &  &  &  &  &  & \textbf{X} &  & \textbf{X} &  &  &  &  & \textbf{X} & \textbf{X} &  &  &  &  &  &  & \textbf{X} & \textbf{X} \\
\textbf{9.2} &  &  &  &  &  &  &  &  &  &  &  &  &  &  &  & \textbf{X} & \textbf{X} &  & \textbf{X} & \textbf{X} &  &  &  &  &  &  & \textbf{X} & \textbf{X} \\
\textbf{9.3} &  & \textbf{X} &  &  &  &  &  &  &  &  &  &  &  & \textbf{X} & \textbf{X} & \textbf{X} &  &  & \textbf{X} & \textbf{X} &  &  &  &  &  &  & \textbf{X} & \textbf{X} \\
\end{tablaSAlarga}

\end{landscape}
\clearpage

\ifproducto
  \section{El reto y el producto final}

\subsection{El reto}
La dirección del IES Ejemplo encarga a cada equipo un informe: \emph{¿dónde se pueden instalar más paneles solares sin que les den sombra los edificios ni las otras filas de paneles?} Para contestar hay que conocer alturas a las que no se llega con una cinta métrica: las del gimnasio y el edificio principal y la del bloque de viviendas del sur, a cuyo pie ni siquiera se puede llegar. El reto está próximo al alumnado: mide su propio instituto y la dirección escucha la respuesta.

\subsection{El modelo matemático}
La figura~\ref{fig:sombra} recoge el modelo en un plano vertical orientado de sur a norte. El mediodía solar es el momento de referencia, y la altura del sol~$h$ se calcula con la latitud~$\varphi$ y la declinación~$\delta$ \parencite{pveducation}:
\[
  h = 90^\circ - \varphi + \delta, \qquad L = \frac{H}{\tan h}.
\]

\begin{figure}[h]
\centering
\begin{tikzpicture}[x=0.42cm,y=0.42cm,>=Stealth]
  \fill[gray!30] (-4,0) rectangle (0,12);
  \draw[thick] (-4,0) rectangle (0,12);
  \draw[thick] (-7,0) -- (28,0);
  \coordinate (T) at (0,12); \coordinate (S) at (24.44,0); \coordinate (O) at (0,0);
  \draw[urjc,thick,->] (-3.5,13.72) -- (T);
  \draw[urjc,thick] (T) -- node[above right,pos=0.45,font=\small]{rayo: $T+t\,\vec v$} (S);
  \draw[very thick,gray!70!black] (O) -- node[below,font=\small]{sombra $L=H/\tan h$} (S);
  \draw[<->] (-5.2,0) -- node[left,font=\small]{$H$} (-5.2,12);
  \draw (20,0) -- (20,2.18);
  \draw[|<->|,font=\footnotesize] (20,-3.2) -- node[below]{$x_P$} (0,-3.2);
  \draw[very thick,blue!60!black] (18.6,0.9) -- (21.4,3.1) node[above right,font=\footnotesize]{panel};
  \fill (T) circle (2.5pt) node[above right,font=\small]{$T=(0,H)$};
  \draw (21.9,0) arc (180:154:2.54);
  \node[font=\small] at (21.2,0.8) {$h$};
  \node[font=\footnotesize,align=center] at (-2,6) {edificio};
  \node[font=\footnotesize] at (-7,1) {sur};
  \node[font=\footnotesize] at (27.5,1) {norte};
\end{tikzpicture}
\caption{Modelo de sombra en el plano vertical sur-norte (elaboración propia).}\label{fig:sombra}
\end{figure}

\begin{itemize}
  \item \textbf{Altura con un punto de medida} (figura~\ref{fig:puntosmedida}a)\textbf{.} Si se puede llegar al pie del objeto, basta medir la distancia $D$ y el ángulo de elevación~$\alpha$: $H = D\tan\alpha + a$, donde $a$ es la altura del ojo. Con $D=24{,}8$~m, $\alpha=25^\circ$ y $a=1{,}6$~m resulta $H\approx 13{,}2$~m, y un grado de error lo cambia a unos $12{,}6$~m o $13{,}7$~m.
  \item \textbf{Altura inaccesible} (figura~\ref{fig:puntosmedida}b)\textbf{.} Si no se puede llegar al pie, como en el bloque de viviendas, se usan dos puntos de medida $A$ y $B$ separados $d$ metros y alineados con el edificio, con ángulos de elevación $\alpha<\beta$. El teorema del seno da
  $H = \dfrac{d\,\sin\alpha\,\sin\beta}{\sin(\beta-\alpha)} + a$. Con $d=10$~m, $\alpha=25^\circ$, $\beta=38^\circ$ y $a=1{,}6$~m resulta $H\approx 13{,}17$~m. Un grado de error en $\alpha$ cambia el resultado a unos $14{,}58$~m o $11{,}95$~m: este método es más sensible que el de un punto de medida, así que solo se usa cuando hace falta y cada medida se repite.
\end{itemize}

\begin{figure}[H]
\centering
\begin{minipage}[b]{0.49\linewidth}\centering
\begin{tikzpicture}[>=Stealth,font=\footnotesize]
  \fill[gray!30] (4.6,0) rectangle (5.4,2.745); \draw[thick] (4.6,0) rectangle (5.4,2.745);
  \draw[thick] (-0.5,0) -- (5.9,0);
  \draw[thick] (0,0) -- (0,0.6); \fill (0,0.6) circle (1.5pt);
  \draw[dashed] (0,0.6) -- (4.6,0.6);
  \draw[urjc,thick] (0,0.6) -- (4.6,2.745); \fill[urjc] (4.6,2.745) circle (1.5pt);
  \draw (1.1,0.6) arc (0:25:1.1); \node at (1.5,0.87) {$\alpha$};
  \draw[<->] (-0.3,0) -- node[left]{$a$} (-0.3,0.6);
  \draw[<->] (0,-0.3) -- node[below]{$D$} (4.6,-0.3);
  \draw[<->] (4.35,0.6) -- node[left]{$D\tan\alpha$} (4.35,2.745);
  \draw[<->] (5.65,0) -- node[right]{$H$} (5.65,2.745);
\end{tikzpicture}
\par\smallskip (a)
\end{minipage}\hfill
\begin{minipage}[b]{0.49\linewidth}\centering
\begin{tikzpicture}[>=Stealth,font=\footnotesize]
  \fill[gray!30] (4.962,0) rectangle (5.7,2.634); \draw[thick] (4.962,0) rectangle (5.7,2.634);
  \fill[urjc!12] (0,0.32) -- (2,0.32) -- (4.962,2.634) -- cycle;
  \draw[thick] (-0.4,0) -- (6.1,0);
  \draw[thick] (0,0) -- (0,0.32) (2,0) -- (2,0.32);
  \fill (0,0.32) circle (1.5pt) node[above left]{$A$}; \fill (2,0.32) circle (1.5pt) node[above left]{$B$};
  \draw[dashed] (0,0.32) -- (4.962,0.32);
  \draw[urjc,thick] (0,0.32) -- (4.962,2.634) (2,0.32) -- (4.962,2.634);
  \fill[urjc] (4.962,2.634) circle (1.5pt);
  \draw (0.95,0.32) arc (0:25:0.95); \node at (1.3,0.5) {$\alpha$};
  \draw (2.6,0.32) arc (0:38:0.6); \node at (2.88,0.58) {$\beta$};
  \node (dif) at (3.2,2.85) {$\beta-\alpha$}; \draw[thin] (dif.east) -- (4.4,2.36);
  \draw[gray!70!black,dashed,thick] (4.3,0) -- (4.3,0.95); \node[below,gray!70!black] at (4.3,0) {valla};
  \draw[<->] (0,-0.3) -- node[below]{$d$} (2,-0.3);
\end{tikzpicture}
\par\smallskip (b)
\end{minipage}
\caption{Altura de un objeto con un punto de medida (a) y con dos puntos de medida cuando no se llega al pie (b) (elaboración propia).}\label{fig:puntosmedida}
\end{figure}

\begin{itemize}
  \item \textbf{Rayo como recta} (figura~\ref{fig:rectaincidencia}a)\textbf{.} El rayo que pasa por $T=(0,H)$ tiene vector director $\vec v=(\cos h,-\sin h)$, es decir, $y = H - x\tan h$. Un punto $P=(x_P,y_P)$ queda en sombra si está por debajo de la recta.
  \item \textbf{Incidencia y ángulo suma} (figura~\ref{fig:rectaincidencia}b)\textbf{.} Con un panel de inclinación $\beta_p$ orientado al sur, de normal $\vec n$ unitaria, y $\vec s$ el vector unitario hacia el sol, el producto escalar da $\cos\theta=\vec n\cdot\vec s=\sin\beta_p\cos h+\cos\beta_p\sin h=\sin(\beta_p+h)$. Aparece así la fórmula del seno de la suma, y la incidencia es perpendicular si $\beta_p + h = 90^\circ$.
\end{itemize}

\begin{figure}[H]
\centering
\begin{minipage}[b]{0.49\linewidth}\centering
\begin{tikzpicture}[>=Stealth,font=\footnotesize]
  \fill[gray!20] (0,0) -- (0,2.6) -- (5.3,0) -- cycle;
  \fill[gray!30] (-0.8,0) rectangle (0,2.6); \draw[thick] (-0.8,0) rectangle (0,2.6);
  \draw[->] (-1,0) -- (5.9,0) node[right]{$x$}; \draw[->] (0,0) -- (0,3.3) node[above]{$y$};
  \draw[urjc,thick] (0,2.6) -- (5.3,0);
  \fill (0,2.6) circle (1.5pt) node[above right]{$T=(0,H)$};
  \draw[->,urjc,very thick] (1.4,1.913) -- ++(1.257,-0.617) node[pos=0.5,above right]{$\vec v$};
  \node[urjc] at (4.2,2.2) {$y=H-x\tan h$};
  \fill (3,0.45) circle (1.5pt) node[left]{$P$ (en sombra)};
  \fill (3.6,1.5) circle (1.5pt) node[right]{$Q$ (al sol)};
  \draw (4.3,0) arc (180:153.85:1); \node at (4.02,0.17) {$h$};
\end{tikzpicture}
\par\smallskip (a)
\end{minipage}\hfill
\begin{minipage}[b]{0.49\linewidth}\centering
\begin{tikzpicture}[>=Stealth,font=\footnotesize]
  \draw[thick] (-0.8,0) -- (3.6,0);
  \node[below] at (-0.5,0) {sur}; \node[below] at (3.3,0) {norte};
  \draw[blue!60!black,very thick] (0,0) -- (2.25,1.3) node[right]{panel};
  \coordinate (M) at (1.125,0.65);
  \draw[dashed] (M) -- (-0.8,0.65);
  \draw[->,thick] (M) -- ++(120:1.6) node[above]{$\vec n$};
  \draw[->,urjc,thick] (M) -- ++(150:2) node[above left]{$\vec s$ (hacia el sol)};
  \draw (M) ++(120:0.85) arc (120:150:0.85); \path (M) ++(135:1.08) node{$\theta$};
  \draw (M) ++(150:1.35) arc (150:180:1.35); \path (M) ++(165:1.58) node{$h$};
  \draw (0.9,0) arc (0:30:0.9); \node at (1.22,0.3) {$\beta_p$};
\end{tikzpicture}
\par\smallskip (b)
\end{minipage}
\caption{El rayo como recta y los puntos en sombra (a), e incidencia del sol sobre un panel orientado al sur (b) (elaboración propia).}\label{fig:rectaincidencia}
\end{figure}

\begin{itemize}
  \item \textbf{Criterio técnico} (figura~\ref{fig:filashoras}a)\textbf{.} El pliego del IDAE fija la distancia mínima $d$ entre una fila y un obstáculo de altura~$h_o$ como $d = h_o / \tan(61^\circ-\text{latitud})$ \parencite{idae}. Para $40{,}4^\circ$ el coeficiente es $k\approx 2{,}66$, es decir, se diseña con un sol de $20{,}6^\circ$, más bajo que los $26{,}15^\circ$ del mediodía del solsticio.
  \item \textbf{Otras horas y orientación} (figura~\ref{fig:filashoras}b)\textbf{.} El plano sur-norte solo describe el mediodía. A otra hora, la altura~$h$ y el ángulo~$A$ que forma el sol con el sur se consultan en la calculadora de la NOAA \parencite{noaa}: la sombra mide $L=H/\tan h$ y su avance hacia el norte es la proyección ortogonal $L\cos A$. Si la fachada no está orientada de este a oeste, el plano se traza hacia el norte desde la arista superior que proyecta la sombra.
\end{itemize}

\begin{figure}[H]
\centering
\begin{minipage}[b]{0.49\linewidth}\centering
\begin{tikzpicture}[>=Stealth,font=\footnotesize]
  \fill[gray!30] (-0.7,0) rectangle (0,1.2); \draw[thick] (-0.7,0) rectangle (0,1.2);
  \draw[thick] (-1.2,0) -- (5,0);
  \draw[gray!70!black,dashed,thick] (0,1.2) -- (2.45,0);
  \draw[urjc,thick] (0,1.2) -- (3.19,0);
  \draw[blue!60!black,very thick] (3.19,0) -- (4.1,0.63) node[right]{fila};
  \draw (2.39,0) arc (180:159.4:0.8);
  \node (ang) at (3.1,0.95) {$61^\circ-\text{latitud}$}; \draw[thin] (ang.south) -- (2.6,0.14);
  \draw[<->] (-0.95,0) -- node[left]{$h_o$} (-0.95,1.2);
  \draw[<->] (0,-0.3) -- node[below]{$d=k\,h_o$} (3.19,-0.3);
\end{tikzpicture}
\par\smallskip (a)
\end{minipage}\hfill
\begin{minipage}[b]{0.49\linewidth}\centering
\begin{tikzpicture}[>=Stealth,font=\footnotesize]
  \fill[gray!30] (-1.3,-0.9) rectangle (1.3,0); \draw[thick] (-1.3,-0.9) rectangle (1.3,0);
  \node at (0,-0.45) {edificio (planta)};
  \draw[dashed] (0,0) -- (0,3.1) node[above]{N};
  \draw[->,very thick,gray!70!black] (0,0) -- node[left]{$L$} (-1.43,2.545);
  \draw[dashed] (-1.43,2.545) -- (0,2.545);
  \draw[<->] (0.35,0) -- node[right]{$L\cos A$} (0.35,2.545);
  \draw (0,0.9) arc (90:119.3:0.9); \path (0,0) ++(104.6:1.2) node{$A$};
  \node[urjc] (sol) at (2.5,-1.05) {sol}; \draw[->,urjc,thick] (sol) -- (1.9,0.02);
\end{tikzpicture}
\par\smallskip (b)
\end{minipage}
\caption{Separación mínima entre filas, con el sol de diseño del IDAE en trazo continuo y el del mediodía en discontinuo (a), y avance de la sombra hacia el norte a otra hora, en planta (b) (elaboración propia).}\label{fig:filashoras}
\end{figure}

Las tablas~\ref{tab:fechas} y~\ref{tab:horas} aplican el modelo a casos concretos. La primera muestra cómo cambia la sombra del mediodía a lo largo del año: el solsticio de invierno es el caso más desfavorable, así que es el que decide la separación. La segunda comprueba, para ese día, que ni siquiera el mediodía basta, porque a otras horas el sol está más bajo.

\begin{table}[H]
\centering
\caption{Altura del sol y sombra de un edificio de 12~m al mediodía solar a 40,4°~N en tres fechas clave.}\label{tab:fechas}
\vspace{4pt}
\begin{tablaSA}{colspec={X[2,l]X[1,c]X[1,c]X[1,c]}, cells={font=\footnotesize}, width=0.8\linewidth}
  Mediodía solar a 40,4° N & Solsticio de invierno & Equinoccio & Solsticio de verano \\
  Declinación $\delta$ & $-23{,}45^\circ$ & $0^\circ$ & $+23{,}45^\circ$ \\
  Altura del sol $h$ & $26{,}15^\circ$ & $49{,}60^\circ$ & $73{,}05^\circ$ \\
  Sombra de un edificio de 12~m & $24{,}44$~m & $10{,}21$~m & $3{,}66$~m \\
\end{tablaSA}
\end{table}

\begin{table}[H]
\centering
\caption{Sombra de un edificio de 12~m en el solsticio de invierno a distintas horas, comparada con el criterio del IDAE.}\label{tab:horas}
\vspace{4pt}
\begin{tablaSA}{colspec={X[2.4,l]X[1,c]X[1,c]X[1,c]X[1,c]}, cells={font=\footnotesize}}
  Solsticio de invierno, edificio de 12~m & Altura del sol & Longitud de sombra & Avance al norte & Coeficiente $k$ \\
  Mediodía solar & $26{,}15^\circ$ & $24{,}44$~m & $24{,}44$~m & $2{,}04$ \\
  Dos horas antes o después ($A\approx 29{,}3^\circ$) & $20{,}31^\circ$ & $32{,}42$~m & $28{,}28$~m & $2{,}36$ \\
  Criterio del IDAE a 40,4° N & $20{,}60^\circ$ & -- & $31{,}93$~m & $2{,}66$ \\
\end{tablaSA}
\end{table}

Como muestra la tabla~\ref{tab:horas}, el mediodía da una separación un 23\,\% menor que la del pliego. Por eso las filas se calculan con el criterio del IDAE, y el modelo explica el motivo: hay que dejar margen para las horas en que el sol está más bajo.

\subsection{Producto final}
Cada equipo entrega un libro de GeoGebra con estas partes:
\begin{enumerate}[leftmargin=1.6em]
  \item Objeto medido, método (uno o dos puntos de medida, con comparación de ambos) y hoja de datos con medidas repetidas.
  \item Construcción dinámica del modelo de sombras con deslizador para la fecha, pensada para que la dirección pueda reutilizarla cambiando la fecha y la altura del obstáculo, y comprobación a mano.
  \item Análisis del error: efecto de un grado de desviación y resultado elegido.
  \item Propuesta de ampliación: plano del patio o cubierta, separación entre filas con el criterio del IDAE, comparada con la del mediodía, y número de paneles que caben, con estimación de producción tomada de PVGIS \parencite{pvgis}.
  \item Reflexión: aportación a los ODS~7 y~11 y límites del modelo (otras horas del día, orientación del edificio, obstáculos no medidos, nubosidad, coste, estructura).
  \item Registro de cambios tras la coevaluación: crítica recibida, qué se modifica y por qué.
  \item Fuentes y autoría de los datos, criterios y construcciones ajenas que se reutilizan.
\end{enumerate}
La propuesta se defiende en cinco minutos ante la clase y la dirección o un representante.

\subsection{Andamiaje}\label{sec:andamiaje}
El apoyo se retira de forma progresiva. Las cuatro fichas, completas en la sección~\ref{sec:fichas}, pasan de una guía paso a paso a un encargo abierto, y cada una prepara lo que la siguiente deja de dar hecho.

\begin{tablaSAlarga}{colspec={Q[l,wd=66pt]X[l]X[l]X[l]X[l]}, cells={font=\scriptsize}, colsep=3pt,
  row{1}={bg=urjc, fg=white, font=\scriptsize\bfseries}, column{1}={bg=rosa, font=\scriptsize\bfseries}, cell{1}{1}={bg=urjc}}
  Aspecto & F1 & F2 & F3 & F4 \\
  Sesión y agrupamiento & S2. Individual, con ayuda del equipo & S5. Individual y parejas & S9. Parejas & S12. Equipo \\
  Nivel de guía & Paso a paso & Guion por objetivos & Encargo con pistas a demanda & Encargo abierto con requisitos \\
  Lugar en la secuencia & Antes de S3 y S4: da significado a las operaciones & Antes de S6: la exposición formaliza lo conjeturado & Tras la exposición de S9: aplica y discrimina casos & Tras S10 y S11: integra los tres bloques \\
  Qué se da hecho & Archivo de partida, datos y orden de los comandos & Lista de elementos de la construcción, sin comandos & Solo el problema y los datos & Solo los requisitos del modelo \\
  Qué decide el alumnado & Qué casos prueba al arrastrar & Cómo construir y qué conjeturar & Cómo construir, qué valores son críticos y qué teorema usar & Cómo definir el rayo y qué ecuación utilizar \\
  Qué debe justificar & Coincidencia entre cálculo a mano y GeoGebra & Una conjetura propia y su demostración & Cuántas soluciones hay y por qué & Una decisión sobre una situación real \\
  Ayuda disponible & Profesor y equipo & Glosario y compañero de pareja & Tres pistas escalonadas, solo si se piden & Tarjetas de ayuda para el equipo, si las pide \\
  Qué se retira & -- & Comandos y archivo de partida & Orden de los pasos & Pistas y vector director \\
  Qué prepara & Rayo y normal en S4 y F4 & Razones de S6 y S7 y sombra de la varilla & Medición con dos puntos de medida (S9 y S14) & Parte~2 del libro (S15) \\
  Criterios & 7.1, 2.1 & 3.2, 3.1 & 1.1, 1.2, 2.1 & 7.2 \\
\end{tablaSAlarga}

Las pistas no restan nota: se registran como información para el docente. Quien no alcanza el nivel~3 de la rúbrica R4 en una ficha recibe un comentario con el paso que debe revisar antes de la ficha siguiente.

El mismo proceso continúa en el producto. Cada sesión magistral termina con una \emph{pieza del reto}, que el equipo guarda para su libro, de modo que S15 y S16 ensamblan lo ya hecho. Se dispone además de plantillas de hoja de datos y de índice del libro, que se usan en S13--S16.

\vspace{4pt}
\begin{tablaSA}{colspec={Q[c,wd=26pt]X[l]X[l]}, cells={font=\footnotesize}}
  Ses. & Pieza del reto & Parte del libro \\
  3 & Vectores entre puntos del plano del centro y distancias reales & 4. Propuesta (plano) \\
  4 & Ángulo entre el rayo solar y la normal de un panel & 4. Propuesta (inclinación) \\
  6 & Altura del sol y sombra de una varilla en varias fechas & 2. Modelo \\
  7 & Inclinación con incidencia perpendicular: $\sin(\beta_p+h)=1$ & 4. Propuesta (inclinación) \\
  8 & Primera medida real: altura del aula con el clinómetro & 1. Método (un punto de medida) \\
  9 & Plan de dos puntos de medida para el bloque de viviendas & 1. Método (dos puntos de medida) \\
  10 & Recta del rayo del edificio elegido y puntos en sombra & 2. Modelo \\
  11 & Filas como rectas paralelas: mediodía frente a IDAE & 4. Propuesta (separación) \\
\end{tablaSA}

\vspace{6pt}
Antes de avanzar hacia el producto, tres hitos comprueban que el equipo está preparado:

\begin{tablaSAlarga}{colspec={X[1.1,l]Q[c,wd=26pt]X[1.6,l]X[1.9,l]X[1.6,l]}, cells={font=\footnotesize},
  row{1}={bg=urjc, fg=white, font=\footnotesize\bfseries}}
  Hito & Ses. & Qué se revisa & Criterio para avanzar & Si no se cumple \\
  Método aprobado & 13 & Objeto, método de medida y reparto de tareas, en unos 3 minutos por equipo & Croquis con las medidas previstas y motivo de usar uno o dos puntos de medida & Se corrige y se reenvía antes de S14 \\
  Datos validados & 14 & Hoja de toma de datos & Tres medidas por magnitud, con unidades, y ángulos que difieren 2° como máximo & Se repite la medida en el cierre de S14 \\
  Borrador coevaluado & 16 & Partes 1 a 4 del libro & Otro equipo aplica R1 y deja al menos una crítica razonada & Se anota en el registro de cambios y se corrige en S17 \\
\end{tablaSAlarga}

\begin{nota}{¿Qué es un clinómetro?}
Instrumento que mide el ángulo de una visual respecto a la horizontal. Se fabrica con un transportador, una pajita como visor y un hilo con plomada; también sirven las aplicaciones de inclinómetro del móvil.
\end{nota}

\subsection{Trabajo cooperativo}\label{sec:cooperativo}
\textbf{Formación de los equipos.} Con el diagnóstico de S1, el profesor forma siete equipos de cuatro estudiantes que combinan distintos niveles y no reúnen a quienes necesitan el mismo apoyo. Si hay ausencias prolongadas y un equipo queda con tres, coordinación y comunicación recaen en la misma persona.

\textbf{Roles y criterios.} El rol no reparte el trabajo: todo el equipo mide, calcula y escribe. Determina quién responde de que cada parte esté bien hecha. Los criterios valen para cualquier pieza, de modo que se pueden observar en todos los turnos.

\begin{tablaSAlarga}{colspec={Q[l,wd=62pt]X[1.2,l]X[2.5,l]X[0.9,l]}, cells={font=\footnotesize},
  row{1}={bg=urjc, fg=white, font=\footnotesize\bfseries}, column{1}={font=\footnotesize\bfseries}}
  Rol & Responsabilidad & Tres criterios de desempeño observables & Evidencia \\
  Coordinación & Organiza las tareas y el tiempo y da la palabra & Mantiene el Kanban actualizado. Asigna a cada tarea responsable y plazo. Comprueba que todos pueden explicar la pieza terminada. & Kanban \\
  Medición & Aporta los datos de cada pieza y garantiza su calidad & Aporta los datos con unidades y origen (plano, tabla o instrumento). Comprueba cada dato con una segunda lectura o medida. Deja registro. En S8 y S14 esto supone calibrar el clinómetro y tomar tres medidas por magnitud. & Pieza, hoja de datos \\
  Modelo & Traslada el problema al cálculo o a GeoGebra & Plantea el cálculo o la construcción. Comprueba el resultado por otra vía. Lo explica al equipo. & Cálculo, construcción \\
  Comunicación & Redacta las piezas y cuida la notación & Usa notación y unidades correctas. Indica de dónde salen datos y fórmulas. Guarda la pieza en su parte del libro; en el turno~4, también el registro de cambios y las publicaciones del TwinSpace. & Libro, TwinSpace \\
\end{tablaSAlarga}

\textbf{Duración y rotación.} Cada estudiante (A, B, C o D en su equipo) ejerce los cuatro roles. El rol solo se ejerce cuando hay trabajo de equipo: S2 y S5 no cuentan porque el trabajo es individual o por parejas, y S18 es el examen.

\begin{tablaSAlarga}{colspec={Q[c,wd=30pt]Q[l,wd=44pt]Q[l,wd=62pt]X[l]Q[l,wd=52pt]Q[l,wd=52pt]Q[l,wd=52pt]Q[l,wd=52pt]}, cells={font=\scriptsize}, colsep=3pt,
  row{1}={bg=urjc, fg=white, font=\scriptsize\bfseries}}
  Turno & Sesiones & Rol activo & Foco & A & B & C & D \\
  1 & S2--S7 & Piezas de S3, S4, S6 y S7 (40 min) & Vectores y trigonometría & Coordinación & Medición & Modelo & Comunicación \\
  2 & S8--S11 & Piezas de S8 a S11 (40 min) & Primera medida real, plan de puntos de medida y rectas & Medición & Modelo & Comunicación & Coordinación \\
  3 & S12--S15 & Cuatro sesiones & F4, mapa, medición y modelo & Modelo & Comunicación & Coordinación & Medición \\
  4 & S16, S17, S19 y S20 & Cuatro sesiones & Propuesta, mejora y defensa & Comunicación & Coordinación & Medición & Modelo \\
\end{tablaSAlarga}

\textbf{Formación en los roles.} Los roles se enseñan; no se da por supuesto que el alumnado sepa ejercerlos. Toda la formación se hace en las sesiones de Matemáticas, una vez formados los equipos.

\begin{tablaSAlarga}{colspec={Q[l,wd=86pt]Q[c,wd=44pt]X[l]}, cells={font=\footnotesize},
  row{1}={bg=urjc, fg=white, font=\footnotesize\bfseries}}
  Momento & Duración & Contenido \\
  S2, inicio & 15 min & Publicación de los equipos, tarjetas de rol del turno~1 y tres normas acordadas por cada equipo. Fórmula de la crítica razonada: «observo\dots, propongo\dots, porque\dots». Cada equipo crea su Kanban (\emph{por hacer}, \emph{en curso}, \emph{hecho}) y asigna la primera tarea. \\
  S3 & Consulta & Página modelo del libro y guía de notación para el rol de comunicación. \\
  S7, S11 y S15 & 3 min & Relevo: quien deja el rol marca sus criterios en la hoja de relevo y anota qué queda pendiente y una mejora; quien lo recibe lee su tarjeta. \\
  S8, pieza del reto & 10 min & Primera medida real: el rol de medición dirige la calibración del clinómetro y el protocolo que se usará en S14. \\
  S16, antes de coevaluar & 3 min & El profesor modela una crítica razonada con un ejemplo de borrador. \\
\end{tablaSAlarga}

\textbf{Interdependencia y responsabilidad individual.} Una pieza no se da por terminada hasta que tiene las aportaciones de los cuatro roles. A la vez, examen, cuestionarios y fichas F1 a F3 son individuales o por parejas, y en la defensa cada miembro expone una parte y cualquiera puede responder a las preguntas, por lo que todos deben conocer el trabajo completo.

\textbf{Evaluación y registro.} Un turno se considera cumplido con al menos dos de los tres criterios del rol, y el indicador de rol de la lista L1 se marca Sí con tres turnos cumplidos de cuatro. Para que sea viable con un solo docente, no se comprueba cada criterio de cada estudiante: al acabar cada turno, cada equipo entrega una hoja de relevo en la que quien ejerció cada rol marca sus criterios y otro miembro lo confirma. Son 28 hojas en toda la secuencia (7 equipos por 4 turnos; la última se completa en S20). El profesor las contrasta con el Kanban y las piezas, observa cada sesión a un equipo distinto y comenta con el equipo cualquier discrepancia. El criterio 9.3 se completa con la coevaluación R3.

\textbf{Conflictos.} El equipo trata primero el problema con la fórmula de la crítica razonada; si no se resuelve, coordinación pide la mediación del profesor, que informa al tutor si hace falta. Solo como último recurso se reorganiza el equipo.

\subsection{Autoría de la propuesta}
Esta situación de aprendizaje es de diseño propio. La medición con clinómetro es una actividad clásica de la enseñanza de la trigonometría; el planteamiento sobre paneles solares es original.

\textbf{ADAPTADA de} \textcite{eratostenes}: la medición simultánea de sombras al mediodía solar en centros distantes del proyecto eTwinning (sección~\ref{sec:etwinning}) se adapta de este proyecto escolar, con finalidad y nivel distintos.

% ======================================================================
\clearpage

\fi
\ifsesiones
  \section{Desarrollo de las sesiones}

\begin{sesionbox}{1}{Inicio}{presentamos el reto}
\campo{Inicio (15 min).} Plano y fotografía de los paneles del gimnasio. Lectura del encargo de la dirección. Relato breve de cómo Eratóstenes midió la Tierra con una sombra. Lluvia de ideas: ¿qué no podemos medir directamente?

\campo{Desarrollo (30 min).} Cuestionario diagnóstico en GeoGebra Classroom \parencite{geogebra}: razones en el triángulo rectángulo, Pitágoras y coordenadas. Presentación del centro socio de eTwinning y de los cuatro roles de equipo.

\campo{Cierre (10 min).} ¿Qué necesitaríamos saber para responder? Cada estudiante estima a ojo la altura del edificio principal y hasta dónde llega su sombra. Preguntas y estimaciones se anotan en un mural que se revisa en S20.

\campo{Agrupamiento:} gran grupo e individual. \campo{Recursos:} proyector, portátiles. \campo{Evaluación:} diagnóstico, sin calificación; sirve para formar equipos heterogéneos y preparar el refuerzo.
\end{sesionbox}

\begin{sesionbox}{2}{Taller}{exploramos los vectores}
\campo{Inicio (15 min).} Se publican los equipos. Formación en los roles: tarjetas del turno~1, normas del equipo, fórmula de la crítica razonada y creación del tablero Kanban con la primera tarea.

\campo{Desarrollo (30 min).} Ficha F1: representar vectores, sumarlos y multiplicarlos por un número arrastrando sus extremos; observar qué ocurre con el producto escalar al cambiar el ángulo y cuándo vale cero. Un vector hace de rayo de sol y otro de normal al panel.

\campo{Cierre (10 min).} Puesta en común: ¿cuándo recibe el panel el sol de frente? Las conjeturas se formalizan en S3 y S4.

\campo{Agrupamiento:} individual con ayuda del equipo. \campo{Recursos:} GeoGebra, F1. \campo{Evaluación:} F1 (7.1, 2.1); observación (9.3).
\end{sesionbox}

\begin{sesionbox}{3}{Magistral}{el conjunto de los vectores}
\campo{Exposición (25 min).} Estructura y propiedades del conjunto de vectores. Suma y producto por un número, a partir de lo observado en F1. Módulo y coordenadas respecto a una base.

\campo{Ejercicios (20 min).} Práctica individual graduada, con corrección comentada.

\campo{Pieza del reto (10 min).} Sobre el plano a escala del centro, cada equipo escribe el vector que va del pie del edificio a una zona candidata y calcula la distancia real.

\campo{Agrupamiento:} gran grupo, individual y equipos. \campo{Recursos:} pizarra, GeoGebra en proyector, hoja de ejercicios, plano del centro. \campo{Evaluación:} seguimiento de los ejercicios, sin calificación.
\end{sesionbox}

\begin{sesionbox}{4}{Magistral}{producto escalar}
\campo{Exposición (25 min).} Definición de producto escalar, $\vec u\cdot\vec v=|\vec u||\vec v|\cos\alpha$, propiedades y expresión en coordenadas. Ángulo entre vectores, proyección ortogonal y perpendicularidad.

\campo{Ejercicios (15 min).} Práctica individual graduada, con corrección comentada.

\campo{Pieza del reto (10 min).} Ejemplo en equipo: ángulo entre el rayo solar del mediodía y la normal de un panel inclinado.

\campo{Cierre (5 min).} Cuestionario Q1 de autocorrección.

\campo{Agrupamiento:} gran grupo, individual y equipos. \campo{Recursos:} pizarra, GeoGebra en proyector, hoja de ejercicios. \campo{Evaluación:} Q1 (7.1, 5.2).
\end{sesionbox}

\begin{sesionbox}{5}{Taller}{la circunferencia goniométrica}
\campo{Inicio (5 min).} Del triángulo rectángulo a un ángulo cualquiera: ¿qué le pasa a la sombra de una varilla cuando sube el sol?

\campo{Desarrollo (40 min).} Ficha F2: construir la circunferencia goniométrica con un deslizador para el ángulo, observar seno, coseno y tangente en los cuatro cuadrantes y formular conjeturas ($\sin^2\alpha+\cos^2\alpha=1$, ángulo doble).

\campo{Cierre (10 min).} Se contrastan las conjeturas; las erróneas se analizan como oportunidad para aprender. Para casa: construir el clinómetro.

\campo{Agrupamiento:} individual y parejas. \campo{Recursos:} GeoGebra, F2. \campo{Evaluación:} F2 (3.2, 3.1); observación (3.2, 9.1).
\end{sesionbox}

\begin{sesionbox}{6}{Magistral}{razones e identidades trigonométricas}
\campo{Exposición (30 min).} Razones de un ángulo cualquiera en grados y radianes. Demostración de la identidad fundamental y de las razones del ángulo suma y diferencia.

\campo{Ejercicios (15 min).} Demostraciones guiadas y cálculo de razones con las fórmulas.

\campo{Pieza del reto (10 min).} Con $h=90^\circ-\varphi+\delta$, cada equipo calcula la altura del sol y la sombra de una varilla de 1~m en varias fechas.

\campo{Agrupamiento:} gran grupo, individual y equipos. \campo{Recursos:} pizarra, hoja de ejercicios, calculadora. \campo{Evaluación:} demostraciones escritas (3.1).
\end{sesionbox}

\begin{sesionbox}{7}{Magistral}{ángulo doble y ecuaciones trigonométricas}
\campo{Exposición (25 min).} Razones del ángulo doble y mitad. Método para resolver ecuaciones trigonométricas sencillas: soluciones en $[0^\circ,360^\circ)$ y solución general. Ejemplo resuelto: $\sin 2x=\cos x$.

\campo{Ejercicios (12 min).} Ecuaciones de dificultad creciente.

\campo{Pieza del reto (10 min).} Resolver $\sin(\beta_p+h)=1$ para hallar la inclinación del panel con incidencia perpendicular en distintas fechas.

\campo{Cierre (8 min).} Cuestionario Q2 (5 min) y relevo de roles (3 min).

\campo{Agrupamiento:} gran grupo, individual y equipos. \campo{Recursos:} pizarra, hoja de ejercicios. \campo{Evaluación:} Q2 (1.2).
\end{sesionbox}

\begin{sesionbox}{8}{Magistral}{teoremas del seno y del coseno}
\campo{Exposición (30 min).} Demostración del teorema del seno y del teorema del coseno. Criterio para elegir cada uno según los datos.

\campo{Ejercicios (15 min).} Resolución de triángulos, justificando qué teorema conviene en cada caso.

\campo{Pieza del reto (10 min).} Primera medida real: con su clinómetro y una cinta, cada equipo mide la altura del aula desde un punto de medida y la compara con el valor real. El rol de medición dirige la calibración del clinómetro.

\campo{Agrupamiento:} gran grupo, individual y equipos. \campo{Recursos:} pizarra, hoja de ejercicios, clinómetro, cinta. \campo{Evaluación:} seguimiento de los ejercicios, sin calificación.
\end{sesionbox}

\begin{sesionbox}{9}{Magistral}{caso ambiguo y altura inaccesible}
\campo{Exposición (20 min).} Caso ambiguo: dos lados y el ángulo opuesto a uno de ellos. Ejemplo resuelto: altura inaccesible desde dos puntos de medida.

\campo{Ficha (20 min).} Ficha F3: comprobar en GeoGebra, arrastrando vértices, cuándo hay dos soluciones, una o ninguna, y justificar cada caso.

\campo{Pieza del reto (10 min).} Cada equipo planifica sobre el plano los dos puntos de medida para medir el bloque de viviendas, al que no se puede llegar.

\campo{Cierre (5 min).} Cuestionario Q3.

\campo{Agrupamiento:} gran grupo, parejas y equipos. \campo{Recursos:} pizarra, GeoGebra, F3, plano del centro. \campo{Evaluación:} F3 (1.1, 1.2, 2.1); Q3 (2.1).
\end{sesionbox}

\begin{sesionbox}{10}{Magistral}{ecuaciones de la recta}
\campo{Exposición (25 min).} Ecuaciones de la recta a partir de un punto y un vector director: vectorial, paramétricas, continua, general y explícita. Vector normal. Posición relativa de puntos y rectas.

\campo{Ejercicios (20 min).} Elegir la ecuación más útil en cada caso y estudiar posiciones relativas.

\campo{Pieza del reto (10 min).} El rayo solar como recta: cada equipo escribe la del edificio elegido y decide si varios puntos del plano quedan en sombra.

\campo{Agrupamiento:} gran grupo, individual y equipos. \campo{Recursos:} pizarra, GeoGebra en proyector, plano del centro. \campo{Evaluación:} seguimiento de los ejercicios, sin calificación.
\end{sesionbox}

\begin{sesionbox}{11}{Magistral}{ángulos y distancias}
\campo{Exposición (25 min).} Incidencia, paralelismo y ortogonalidad. Ángulo entre rectas y distancia de un punto a una recta mediante el producto escalar.

\campo{Ejercicios (12 min).} Problemas métricos de dificultad creciente.

\campo{Pieza del reto (10 min).} Las filas de paneles como rectas paralelas: cada equipo compara el coeficiente del mediodía con el del IDAE y explica por qué el pliego exige más separación.

\campo{Cierre (8 min).} Cuestionario Q4 (5 min) y relevo de roles (3 min).

\campo{Agrupamiento:} gran grupo, individual y equipos. \campo{Recursos:} pizarra, GeoGebra en proyector, hoja de ejercicios. \campo{Evaluación:} Q4 (5.2).
\end{sesionbox}

\begin{sesionbox}{12}{Taller}{rayo y sombra}
\campo{Inicio (5 min).} Cada estudiante revisa la tarjeta de su nuevo rol y se recuerda qué pide el encargo.

\campo{Desarrollo (40 min).} Ficha F4 en equipo. Todos los miembros trabajan juntos; el equipo puede pedir tarjetas de ayuda, y quien termina antes aborda la ampliación.

\campo{Cierre (10 min).} Puesta en común sobre qué ecuación conviene (5 min). Cada equipo completa el croquis de su método a partir de la pieza de S9 (5 min).

\campo{Agrupamiento:} equipos. \campo{Recursos:} GeoGebra, F4, tarjetas de ayuda. \campo{Evaluación:} F4 (7.2); observación (9.1).
\end{sesionbox}

\begin{sesionbox}{13}{Taller}{refuerzo, mapa conceptual y método}
\campo{Inicio (5 min).} Devolución de los cuestionarios y reparto por criterios.

\campo{Refuerzo y ampliación (20 min).} Trabajo individual o por parejas. Quien lo necesita revisa, con ejercicios resueltos paso a paso, el criterio que indiquen sus cuestionarios: 7.1 y 5.2 (Q1), 1.2 (Q2), 2.1 (Q3) o 5.2 (Q4). El resto aborda la ampliación: panel inclinado óptimo o más de un obstáculo en el modelo.

\campo{Mapa y método (25 min).} Cada equipo elabora el mapa conceptual que relaciona vectores, trigonometría y rectas, a partir de una lista de conceptos. Mientras, el profesor pasa por cada equipo, unos 3 minutos, para aprobar su método de medida.

\campo{Cierre (5 min).} Revisión del Kanban y preparación de la salida.

\campo{Agrupamiento:} individual, parejas y equipos. \campo{Recursos:} ejercicios resueltos, plantilla del mapa, croquis. \campo{Evaluación:} mapa (5.1); segundo intento de los cuestionarios antes de S18.
\end{sesionbox}

\begin{sesionbox}{14}{Tarea}{medimos el instituto}
\campo{Preparación.} La sesión se sitúa en la franja que incluye el mediodía solar (hacia las 13:25 en la segunda semana de marzo; la hora exacta se consulta en la calculadora de la NOAA \parencite{noaa}).

\campo{Desarrollo (45 min).} En el patio: medir con clinómetro y cinta, repitiendo cada medida. Se mide un objeto accesible desde uno y desde dos puntos de medida, para comparar los métodos, y el bloque de viviendas desde dos puntos de medida. Medir la sombra de una varilla vertical de 1~m al mediodía solar y subir el dato al TwinSpace.

\campo{Si hay nubes.} Las alturas se miden igual, porque el clinómetro no necesita sol. La sombra de la varilla la mide un representante de cada equipo el primer día despejado de esa semana, y se anota la fecha. Si no hay ninguno, se usa la altura del sol que da la NOAA y se indica en el libro.

\campo{Cierre (10 min).} Cada equipo valida sus datos y detecta medidas que debe repetir.

\campo{Agrupamiento:} equipos con roles. \campo{Recursos:} clinómetro, cinta, varilla, hoja de datos. \campo{Evaluación:} observación (9.1, 9.3).
\end{sesionbox}

\begin{sesionbox}{15}{Tarea}{modelo de sombras y análisis del error}
\campo{Inicio (5 min).} Objetivos del día en el Kanban.

\campo{Desarrollo (40 min).} A partir de las piezas del reto ya resueltas: modelo a escala en GeoGebra con deslizador para la fecha, resolución de triángulos y comprobación a mano. Efecto de un grado de error y elección del resultado. Sombras en solsticios y equinoccios.

\campo{Cierre (10 min).} Prueba de reutilización: cambiar fecha y altura y comprobar que el modelo responde. Actualización del libro y relevo de roles (3 min).

\campo{Agrupamiento:} equipos. \campo{Recursos:} GeoGebra, piezas del reto, hoja de datos. \campo{Evaluación:} libro (1.1, 2.2, 7.2); observación (9.3).
\end{sesionbox}

\begin{sesionbox}{16}{Tarea}{propuesta de ampliación y coevaluación}
\campo{Inicio (5 min).} Objetivos del día en el Kanban.

\campo{Desarrollo (30 min).} Cálculo de filas con el criterio del IDAE, plano de la propuesta, número de paneles y producción con PVGIS, con la orientación del departamento de Tecnología. Comparación con los datos del socio y reflexión sobre ODS y límites del modelo.

\campo{Cierre (20 min).} El profesor modela una crítica razonada (3 min). Coevaluación del borrador entre equipos con R1 (12 min) y anotación de las críticas en el registro de cambios (5 min).

\campo{Agrupamiento:} equipos. \campo{Recursos:} GeoGebra, PVGIS, TwinSpace, R1. \campo{Evaluación:} libro y propuesta (6.1, 6.2, 7.2); coevaluación (9.2).
\end{sesionbox}

\begin{sesionbox}{17}{Tarea}{mejora del libro y ensayo}
\campo{Inicio (5 min).} Cada equipo prioriza las críticas recibidas.

\campo{Desarrollo (30 min).} Se incorporan las críticas, se completan las partes pendientes y se cierra el registro de cambios. Cada miembro prepara la parte que va a exponer.

\campo{Ensayo (20 min).} Cada equipo ensaya ante otro, con tiempo controlado, y recibe una mejora concreta.

\campo{Agrupamiento:} equipos. \campo{Recursos:} libros, R2. \campo{Evaluación:} registro de cambios (9.2); observación (9.3). El libro se entrega al terminar la sesión.
\end{sesionbox}

\begin{sesionbox}{18}{Evaluación}{examen y retroalimentación}
\campo{Examen (45 min).} Prueba individual escrita (sección~\ref{sec:examen}).

\campo{Cierre (10 min).} Autocorrección guiada de un ejercicio; la devolución individual se entrega antes de S20.

\campo{Agrupamiento:} individual. \campo{Recursos:} examen, calculadora. \campo{Evaluación:} examen (1.1, 1.2, 2.1, 3.1, 5.1, 5.2, 7.2).
\end{sesionbox}

\begin{sesionbox}{19}{Cierre}{defendemos la propuesta (I)}
\campo{Apertura (3 min).} Orden de intervención y equipos asignados para coevaluar.

\campo{Defensas (32 min).} Cuatro equipos ante la clase y la dirección o un representante: 5 minutos de exposición, 2 de preguntas y 1 de cambio.

\campo{Coevaluación (15 min).} Cada equipo asignado completa la rúbrica del que ha valorado.

\campo{Cierre (5 min).} Primeras conclusiones sobre las propuestas.

\campo{Agrupamiento:} equipos y gran grupo. \campo{Recursos:} libros, proyector, rúbricas. \campo{Evaluación:} exposición (1.1, 2.2, 6.1, 6.2, 8.1); coevaluación (7.2, 8.1).
\end{sesionbox}

\begin{sesionbox}{20}{Cierre}{defendemos la propuesta (II) y cerramos}
\campo{Defensas (24 min).} Los tres equipos restantes, con el mismo formato.

\campo{Coevaluación (10 min).} Rúbrica de los equipos asignados y valoración dentro del propio equipo (9.1, 9.2, 9.3).

\campo{Metacognición (13 min).} Se retoman el mural de S1 y las estimaciones iniciales; cada estudiante completa la autoevaluación y cada equipo, la hoja de relevo del turno~4.

\campo{Cierre (8 min).} Videomensaje conjunto para el socio eTwinning.

\campo{Agrupamiento:} equipos y gran grupo. \campo{Recursos:} libros, proyector, rúbricas. \campo{Evaluación:} exposición (1.1, 2.2, 6.1, 6.2, 8.1); co/autoevaluación (7.2, 8.1, 9.1, 9.2, 9.3).
\end{sesionbox}

\fi
\section{Evaluación}

\subsection{Planteamiento}
La evaluación es continua y sirve para mejorar la enseñanza y el aprendizaje \parencite[art.~20.1, p.~22]{decreto64}. Tiene tres momentos: diagnóstica (S1), formativa (cuestionarios, fichas y coevaluación del borrador) y sumativa (examen, libro y defensa). Participan el profesor, los iguales y el propio alumnado.

\subsection{Criterios de calificación}\label{sec:calificacion}
\begin{tablaSA}{colspec={X[2.3,l]Q[c,wd=30pt]X[1.7,l]X[2.4,l]}, cells={font=\footnotesize}}
  Instrumento & Peso & Criterios & Nota del instrumento (0--10) \\
  Examen escrito & 35\,\% & 1.1, 1.2, 2.1, 3.1, 5.1, 5.2, 7.2 & Suma de los puntos del examen \\
  Libro de GeoGebra y propuesta (R1) & 25\,\% & 1.1, 2.2, 6.1, 6.2, 7.2 & Media ponderada de los niveles de R1 \\
  Fichas de GeoGebra F1--F4 (R4) & 10\,\% & 1.1, 1.2, 2.1, 3.1, 3.2, 7.1, 7.2 & Media de las cuatro fichas, cada una valorada con R4 \\
  Exposición oral (R2) & 10\,\% & 1.1, 2.2, 6.1, 6.2, 8.1 & Media ponderada de los niveles de R2 \\
  Cuestionarios Q1--Q4 & 5\,\% & 1.2, 2.1, 5.2, 7.1 & Media de los cuatro: porcentaje de aciertos sobre 10, con el mejor de dos intentos \\
  Mapa conceptual (L2) & 5\,\% & 5.1 & (n.º de Sí / 3) $\times$ 10 \\
  Observación sistemática (L1) & 5\,\% & 3.2, 9.1, 9.2, 9.3 & (n.º de Sí / 6) $\times$ 10 \\
  Coevaluación y autoevaluación (R3) & 5\,\% & 7.2, 8.1, 9.1, 9.2, 9.3 & Media de las casillas auto y co, $\times$ 2,5 \\
\end{tablaSA}

La nota del libro es común al equipo, ajustada individualmente si la observación o la coevaluación muestran diferencias claras de participación.

\subsection{Retroalimentación y refuerzo}
\begin{itemize}
  \item Los cuestionarios se autocorrigen y determinan qué criterio refuerza cada estudiante en S13; el segundo intento se hace después de ese refuerzo.
  \item Las fichas se devuelven con un comentario sobre qué mejorar.
  \item El borrador del libro se revisa con la rúbrica antes de su entrega, de modo que la crítica pueda incorporarse.
  \item Tras el examen, quien no alcance el nivel adecuado en algún criterio recibe tareas específicas y una nueva oportunidad en la siguiente unidad.
\end{itemize}

\subsection{Instrumentos}\label{sec:instrumentos}

\subsubsection{R1. Rúbrica del libro de GeoGebra y la propuesta}
\begin{tablaSAlarga}{colspec={Q[l,wd=66pt]X[l]X[l]X[l]X[l]}, cells={font=\scriptsize}, colsep=3pt,
  row{1}={bg=urjc, fg=white, font=\scriptsize\bfseries}, column{1}={bg=rosa, font=\scriptsize\bfseries}, cell{1}{1}={bg=urjc}}
  Criterio (peso) & 1 · Inicial & 2 · En desarrollo & 3 · Adecuado & 4 · Excelente \\
  Modelo y método, 1.1 (30\,\%) & Sin modelo a escala o incorrecto & Modelo incompleto; método poco justificado & Modelo correcto, reutilizable con otra fecha y altura, y método explicado & Además compara uno y dos puntos de medida y valora su eficiencia \\
  Error y solución, 2.2 (20\,\%) & Una sola medida sin análisis & Repite medidas, sin análisis del error & Analiza un grado de error y elige un resultado & Además justifica la decisión según el contexto \\
  Representación, 7.2 (20\,\%) & Representaciones confusas & Claras pero sin relacionar & Integra construcción, cálculos y explicación, y cita sus fuentes & Además elige y justifica las representaciones más útiles \\
  Propuesta, 6.1 (20\,\%) & Sin cálculo de filas ni plano & Plano sin aplicar el criterio técnico & Plano y filas correctos con el criterio del IDAE & Además estima producción y plantea alternativas \\
  ODS y límites, 6.2 (10\,\%) & No se menciona & Mención general & Relaciona con los ODS y señala un límite & Analiza el impacto y varios límites del modelo \\
\end{tablaSAlarga}

\subsubsection{R2. Rúbrica de la exposición oral}
\vspace*{10pt}
\begin{tablaSAlarga}{colspec={Q[l,wd=66pt]X[l]X[l]X[l]X[l]}, cells={font=\scriptsize}, colsep=3pt,
  row{1}={bg=urjc, fg=white, font=\scriptsize\bfseries}, column{1}={bg=rosa, font=\scriptsize\bfseries}, cell{1}{1}={bg=urjc}}
  Criterio (peso) & 1 · Inicial & 2 · En desarrollo & 3 · Adecuado & 4 · Excelente \\
  Organización y rigor, 8.1 (40\,\%) & Desordenada, sin notación & Orden parcial, errores de notación & Ordenada, terminología correcta & Además adaptada al público y sin errores \\
  Método y solución, 1.1 y 2.2 (30\,\%) & No explica el método & Lo explica sin justificar & Explica y justifica el resultado & Además responde con solvencia a preguntas \\
  Propuesta, 6.1 y 6.2 (30\,\%) & No formula propuesta & Propuesta sin datos & Propuesta con datos y ODS & Además argumenta impacto y viabilidad \\
\end{tablaSAlarga}

\subsubsection{R3. Coevaluación y autoevaluación}
Escala de 1 (nunca) a 4 (siempre). La autoevaluación la completa cada estudiante en S20. En la coevaluación, un equipo asignado por el profesor valora los indicadores 7.2 y 8.1, y los compañeros del propio equipo, los indicadores 9.1, 9.2 y 9.3; si hay varias valoraciones, se hace la media. Si auto y co difieren en 2 puntos o más, el profesor lo comenta con el estudiante antes de cerrar la nota.

\begin{tablaSA}{colspec={X[l]Q[c,wd=30pt]Q[c,wd=62pt]Q[c,wd=44pt]Q[c,wd=44pt]}, cells={font=\footnotesize}}
  Indicador & Crit. & Coevalúa & Auto (1--4) & Co (1--4) \\
  Las representaciones del libro permiten entender el método & 7.2 & Otro equipo & & \\
  La explicación se sigue con orden y buena notación & 8.1 & Otro equipo & & \\
  Uso mis errores y los del equipo para revisar el razonamiento & 9.1 & Propio equipo & & \\
  Incorporo la crítica razonada para mejorar & 9.2 & Propio equipo & & \\
  Asumo mi rol y escucho el razonamiento de los demás & 9.3 & Propio equipo & & \\
\end{tablaSA}

\pista{Metacognición: ¿qué sabía antes? ¿Qué me costó más? ¿Qué haría distinto en la próxima medición?}

\subsubsection{L1. Lista de cotejo de observación y Kanban}
Salvo el indicador de rol, que tiene su propia regla, un indicador se marca Sí si se observa al menos una vez en las sesiones que la sección~\ref{sec:matrices} asigna a su criterio.

\begin{tablaSA}{colspec={X[l]Q[c,wd=30pt]Q[c,wd=30pt]Q[c,wd=30pt]}, cells={font=\footnotesize}}
  Indicador (por estudiante) & Crit. & Sí & No \\
  Investiga en GeoGebra antes de pedir ayuda & 3.2 & & \\
  Analiza un error propio y lo corrige & 9.1 & & \\
  Persevera ante una medida o construcción fallida & 9.2 & & \\
  Aporta al registro de cambios del libro la crítica recibida y lo que modifica & 9.2 & & \\
  Cumple los criterios de su rol\soloProducto{ (sección~\ref{sec:cooperativo})} en al menos tres de los cuatro turnos & 9.3 & & \\
  Escucha y tiene en cuenta el razonamiento ajeno & 9.3 & & \\
\end{tablaSA}

\subsubsection{L2. Lista de cotejo del mapa conceptual}
\begin{tablaSA}{colspec={X[l]Q[c,wd=30pt]Q[c,wd=30pt]Q[c,wd=30pt]}, cells={font=\footnotesize}}
  Indicador & Crit. & Sí & No \\
  Relaciona vector director y normal con las ecuaciones de la recta & 5.1 & & \\
  Conecta producto escalar con ángulos y distancias & 5.1 & & \\
  Indica qué herramienta usa para cada tipo de problema & 5.1 & & \\
\end{tablaSA}

\subsubsection{R4. Rúbrica común de las fichas de GeoGebra}
\begin{tablaSAlarga}{colspec={Q[l,wd=66pt]X[l]X[l]X[l]X[l]}, cells={font=\scriptsize}, colsep=3pt,
  row{1}={bg=urjc, fg=white, font=\scriptsize\bfseries}, column{1}={bg=rosa, font=\scriptsize\bfseries}, cell{1}{1}={bg=urjc}}
  Aspecto (peso) & 1 · Inicial & 2 · En desarrollo & 3 · Adecuado & 4 · Excelente \\
  Construcción (30\,\%) & No funciona & Funciona con errores o incompleta & Correcta y dinámica & Además está etiquetada y se puede reutilizar \\
  Resultado y comprobación (30\,\%) & Sin resultado & Resultado sin comprobar & Resultado correcto y comprobado a mano & Además explica las discrepancias \\
  Respuesta justificada (40\,\%) & No responde & Responde sin justificar & Justifica con un razonamiento correcto & Además generaliza o propone otro caso \\
\end{tablaSAlarga}

\pista{Criterios que recoge cada aspecto. Construcción: 7.1 en F1, 3.2 en F2 y 7.2 en F4. Resultado y comprobación: 2.1 en F1 y F3, y 7.2 en F4. Respuesta justificada: 2.1 en F1, 3.1 en F2, 1.1 y 1.2 en F3, y 7.2 en F4.}

\subsubsection{Fichas de GeoGebra}\label{sec:fichas}
\soloProducto{Las cuatro fichas siguen la progresión descrita en la sección~\ref{sec:andamiaje}. }Se valoran con la rúbrica R4.

\begin{tarjeta}[Paso a paso · sesión 2 · criterios 7.1 y 2.1]{F1. Vectores y producto escalar}
\begin{enumerate}[leftmargin=1.5em]
  \item Abre el archivo de partida, con ejes y cuadrícula. Crea $\vec u=(4,3)$ y $\vec v=(1,2)$ con la herramienta Vector.
  \item Crea $\vec u+\vec v$ y $2\vec u$. Arrastra los extremos y describe qué figura forman $\vec u$, $\vec v$ y $\vec u+\vec v$.
  \item Muestra $\vec u\cdot\vec v$ y el ángulo que forman. Arrastra $\vec v$ hasta que el producto escalar valga 0.
  \item Vuelve a $\vec v=(1,2)$. Calcula a mano $|\vec u|$, el ángulo y la proyección de $\vec u$ sobre $\vec v$, y compáralos con GeoGebra.
  \item Sea $\vec s$ el vector hacia el sol y $\vec n$ la normal al panel. ¿Cuándo es máximo $\vec s\cdot\vec n$? ¿Qué significa que valga 0?
\end{enumerate}
\campo{Pregunta final.} ¿Cuándo es nulo el producto escalar? Compruébalo a mano con un ejemplo.

\campo{Solución.} $\vec u\cdot\vec v=10$, $|\vec u|=5$, $|\vec v|=\sqrt5$, $\cos\alpha=2/\sqrt5$, $\alpha\approx26{,}57^\circ$ y proyección $10/\sqrt5\approx4{,}47$. Dos vectores no nulos tienen producto escalar nulo si son perpendiculares: $(4,3)\cdot(-3,4)=0$. El máximo se da con el sol de frente; el valor 0 corresponde a un rayo rasante.

\campo{Refuerzo:} captura de cada paso. \campo{Ampliación:} halla los vectores unitarios perpendiculares a $\vec u$.
\end{tarjeta}

\begin{tarjeta}[Guion por objetivos · sesión 5 · criterios 3.2 y 3.1]{F2. La circunferencia goniométrica}
Tu construcción debe incluir:
\begin{enumerate}[leftmargin=1.5em]
  \item Un deslizador $\alpha\in[0^\circ,360^\circ]$, la circunferencia de radio 1 y un punto $P$ que la recorra.
  \item Segmentos que muestren el seno y el coseno de $\alpha$, y el valor de la tangente.
  \item Una tabla con los signos de las tres razones en cada cuadrante.
  \item El valor de $\sin^2\alpha+\cos^2\alpha$ mientras se mueve el deslizador.
  \item El cociente $\sin 2\alpha/(\sin\alpha\cos\alpha)$. Un compañero propone que $\sin 2\alpha=2\sin\alpha$: busca un contraejemplo y propón una relación correcta.
\end{enumerate}
\campo{Pregunta final.} Conjetura una identidad y demuéstrala por escrito (la demostración se completa en S6).

\campo{Solución.} $\sin^2\alpha+\cos^2\alpha=1$: es el teorema de Pitágoras en un triángulo de hipotenusa 1, válido en todos los cuadrantes porque las razones aparecen al cuadrado. Contraejemplo: con $\alpha=30^\circ$, $\sin60^\circ\approx0{,}866\neq1$. Donde el denominador no se anula, el cociente vale 2, luego $\sin2\alpha=2\sin\alpha\cos\alpha$, que se demuestra en S7 haciendo $\beta=\alpha$ en el ángulo suma.

\campo{Refuerzo:} se da la definición $P=(\cos\alpha,\sin\alpha)$. \campo{Ampliación:} conjetura una expresión para $\cos2\alpha$.
\end{tarjeta}

\begin{tarjeta}[Encargo con pistas · sesión 9 · criterios 1.1, 1.2 y 2.1]{F3. ¿Cuántos triángulos hay?}
\campo{Encargo.} Construye un triángulo con $A=30^\circ$, $b=10$ y un deslizador para el lado $a$, de modo que se vean todas las posiciones posibles de $B$.
\begin{enumerate}[leftmargin=1.5em]
  \item ¿Para qué valores de $a$ hay dos triángulos, uno o ninguno? Justifícalo con un cálculo.
  \item Para $a=8$, resuelve a mano todos los triángulos y compruébalos arrastrando el deslizador. Si falta tiempo, este apartado se termina en casa.
  \item Si conocieras dos lados y el ángulo comprendido, ¿qué teorema usarías? ¿Y en este caso? Razónalo.
\end{enumerate}
\campo{Pistas, solo si se piden.} 1: ¿qué figura forman los puntos situados a distancia $a$ de $C$? 2: ¿a qué distancia está $C$ de la recta $AB$? 3: compara $a$ con $b\sin A$ y con $b$.

\campo{Solución.} $b\sin A=5$. Ninguno si $a<5$; uno, rectángulo, si $a=5$; dos si $5<a<10$; uno si $a\ge10$. Para $a=8$, $\sin B=0{,}625$: $B\approx38{,}68^\circ$, $C\approx111{,}32^\circ$, $c\approx14{,}91$; o bien $B\approx141{,}32^\circ$, $C\approx8{,}68^\circ$, $c\approx2{,}42$. Con dos lados y el ángulo comprendido conviene el teorema del coseno; aquí, con un lado y su ángulo opuesto, el del seno.

\campo{Refuerzo:} las pistas 1 y 2 se dan desde el inicio. \campo{Ampliación:} estudia qué ocurre si $A$ es obtuso.
\end{tarjeta}

\begin{tarjeta}[Encargo abierto · sesión 12 · criterio 7.2]{F4. Rayo y sombra}
\campo{Encargo.} Construye un modelo dinámico de la sombra de un edificio que cumpla estos requisitos:
\begin{enumerate}[leftmargin=1.5em]
  \item Deslizadores $h\in[5^\circ,85^\circ]$ y $H\in[0,20]$, y el punto $T=(0,H)$.
  \item El rayo que pasa por $T$, definido con un vector director que debes deducir a partir de $h$.
  \item El punto $S$ donde el rayo corta al eje $OX$, comprobando que su abscisa es $H/\tan h$.
  \item Un panel en el punto $P=(20,3)$, con $H=12$.
  \item La recta en forma explícita y general. ¿Cuál te sirve mejor para decidir si $P$ queda en sombra? Justifícalo.
\end{enumerate}
\campo{Pregunta final.} ¿Para qué valores de $h$ queda $P$ en sombra?

\campo{Solución.} $P$ está en sombra si $3<12-20\tan h$, es decir, si $\tan h<0{,}45$, luego $h<24{,}23^\circ$. Un vector director es $\vec v=(\cos h,-\sin h)$. La forma explícita es la más cómoda, porque permite comparar directamente la altura de la recta con la de $P$.

\campo{Tarjetas de ayuda para el equipo, si las pide:} el vector $\vec v$, los pasos con capturas y datos prefijados. \campo{Ampliación:} añade un panel inclinado y calcula $h$ para incidencia perpendicular.
\end{tarjeta}

\subsubsection{Hoja de toma de datos}
\begin{tablaSA}{colspec={X[2,l]X[1,c]X[1,c]X[1,c]X[1,c]}, cells={font=\footnotesize}}
  Magnitud (con unidad) & Medida 1 & Medida 2 & Medida 3 & Valor elegido \\
  Altura del ojo $a$ (m) & & & & \\
  Un punto de medida: distancia al pie $D$ (m) & & & & \\
  Un punto de medida: ángulo de elevación (°) & & & & \\
  Dos puntos de medida: distancia entre ellos $d$ (m) & & & & \\
  Dos puntos de medida: ángulo $\alpha$ en A (°) & & & & \\
  Dos puntos de medida: ángulo $\beta$ en B (°) & & & & \\
  Sombra de la varilla de 1~m (m), con fecha y hora & & & & \\
\end{tablaSA}

\subsubsection{Cuestionario modelo}
Banco de preguntas de autocorrección en GeoGebra Classroom, con respuesta entre paréntesis.
\begin{enumerate}[label=\textbf{\arabic*.},leftmargin=1.6em]
  \item (Q1, 7.1) Si $\vec u=(1,2)$ y $\vec v=(3,-1)$, ¿cuánto vale $\vec u\cdot\vec v$? \emph{(1)}
  \item (Q1, 5.2) ¿Para qué $k$ son perpendiculares $(2,k)$ y $(3,6)$? \emph{($k=-1$)}
  \item (Q2, 1.2) Soluciones de $2\sin x=1$ en $[0^\circ,360^\circ)$. \emph{($30^\circ$ y $150^\circ$)}
  \item (Q3, 2.1) Con $a=3$, $b=5$ y $A=80^\circ$, ¿existe el triángulo? \emph{(No: $\sin B\approx 1{,}64>1$)}
  \item (Q4, 5.2) Distancia de $(0,0)$ a $3x+4y-10=0$. \emph{(2)}
\end{enumerate}

\subsubsection{Examen modelo}\label{sec:examen}
\begin{tarjeta}[Individual · 45 minutos · calculadora permitida]{Examen. Medir lo que no se puede alcanzar}
\begin{enumerate}[label=\textbf{\arabic*.},leftmargin=1.6em,itemsep=3pt]
  \item \textbf{(2 puntos; 5.2)} Dados $\vec u=(3,1)$ y $\vec v=(2,-1)$: a) calcula $\vec u+2\vec v$ y $\vec u\cdot\vec v$; b) halla el ángulo que forman y la proyección de $\vec u$ sobre $\vec v$; c) halla $k$ para que $\vec w=(k,2)$ sea perpendicular a $\vec u$.
  \item \textbf{(1 punto; 2.1)} Un compañero afirma que la proyección de $\vec u$ sobre $\vec v$ vale $5/\sqrt{10}$, pero GeoGebra muestra $2{,}236$. ¿Quién tiene razón? Explica el origen del error.
  \item \textbf{(2 puntos; 3.1, 1.2)} a) A partir de la fórmula de $\sin(\alpha+\beta)$, demuestra que $\sin 2\alpha=2\sin\alpha\cos\alpha$. b) Resuelve $\sin 2x=\cos x$ dando todas las soluciones.
  \item \textbf{(2,5 puntos; 1.1, 1.2)} a) Desde $A$ se ve lo alto de un edificio con un ángulo de $25^\circ$ y, acercándose 10~m, con $38^\circ$. El ojo está a $1{,}6$~m del suelo. Calcula la altura y justifica qué teorema has elegido. b) Resuelve el triángulo con $a=8$, $b=10$ y $A=40^\circ$.
  \item \textbf{(2,5 puntos; 5.1, 5.2, 7.2)} a) Escribe las ecuaciones de la recta $r$ que pasa por $A(1,2)$ con vector director $(2,-1)$. b) Estudia su posición respecto a $s: 3x-4y+5=0$ y halla el ángulo que forman. c) Calcula la distancia de $P(4,3)$ a $r$. d) Un edificio de 12~m tiene su fachada norte en $x=0$ y el sol está a $26{,}15^\circ$: ¿queda en sombra un panel en $(20,3)$? Indica qué ecuación de la recta has preferido y por qué.
\end{enumerate}
\end{tarjeta}

\begin{tarjeta}{Solución del examen}
\small
\textbf{1.} a) $(7,-1)$; $\vec u\cdot\vec v=5$. b) $\cos\alpha=5/(\sqrt{10}\sqrt5)=\sqrt2/2$, luego $\alpha=45^\circ$; proyección $5/\sqrt5=\sqrt5$. c) $3k+2=0$, $k=-2/3$.

\textbf{2.} GeoGebra: $\sqrt5\approx2{,}236$. El compañero dividió por $|\vec u|$, con lo que calculó la proyección de $\vec v$ sobre $\vec u$.

\textbf{3.} a) Tomando $\beta=\alpha$. b) $\cos x(2\sin x-1)=0$: $x=90^\circ+180^\circ k$, $x=30^\circ+360^\circ k$, $x=150^\circ+360^\circ k$, con $k\in\mathbb Z$.

\textbf{4.} a) En el triángulo se conocen un lado y dos ángulos, por lo que se usa el teorema del seno: $H=\dfrac{10\sin25^\circ\sin38^\circ}{\sin13^\circ}+1{,}6\approx13{,}17$~m. b) Dos soluciones: $B\approx53{,}46^\circ$, $C\approx86{,}54^\circ$, $c\approx12{,}42$; o bien $B\approx126{,}54^\circ$, $C\approx13{,}46^\circ$, $c\approx2{,}90$.

\textbf{5.} a) $(x,y)=(1,2)+t(2,-1)$; $x=1+2t$, $y=2-t$; $\frac{x-1}{2}=\frac{y-2}{-1}$; $x+2y-5=0$; $y=-\frac12x+\frac52$. b) Secantes, pues $(1,2)$ y $(3,-4)$ no son proporcionales; $\cos\theta=\frac{|3-8|}{\sqrt5\cdot5}=\frac1{\sqrt5}$, $\theta\approx63{,}43^\circ$. c) $d=\frac{|4+6-5|}{\sqrt5}=\sqrt5$. d) En $x=20$ la recta del rayo tiene altura $12-20\tan26{,}15^\circ\approx2{,}18<3$: el panel no queda en sombra; la forma explícita permite comparar alturas directamente.
\end{tarjeta}

\section{Atención a la diversidad}\label{sec:diversidad}
El decreto regula las medidas de apoyo y las condiciones de accesibilidad y diseño universal, y se rige por los principios de normalización e inclusión \parencite[art.~32, pp.~27--28]{decreto64}. La propuesta sigue las pautas del Diseño Universal para el Aprendizaje \parencite{cast2024}.

\begin{tablaSAlarga}{colspec={Q[l,wd=100pt]X[l]}, row{1}={bg=urjc, fg=white, font=\footnotesize\bfseries}, column{1}={font=\footnotesize\bfseries}}
  Principio DUA & Medidas para todo el grupo \\
  Implicación & Reto real con destinatario; elección del objeto a medir; roles rotatorios; error tratado como oportunidad. \\
  Representación & Explicación magistral, construcción dinámica y material escrito; glosario de geometría disponible durante toda la secuencia; ejemplos con el propio centro. \\
  Acción y expresión & Productos variados (libro, exposición, escrito, mapa conceptual); clinómetro casero o digital; plantillas de planificación y Kanban. \\
\end{tablaSAlarga}

\begin{tablaSAlarga}{colspec={Q[l,wd=100pt]X[2,l]X[1,l]}, row{1}={bg=urjc, fg=white, font=\footnotesize\bfseries}, column{1}={font=\footnotesize\bfseries}}
  Necesidad & Medidas y actividades & Elementos afectados \\
  Refuerzo (según Q1--Q4) & En S13, ejercicios resueltos paso a paso y parejas de apoyo sobre el criterio que indiquen los cuestionarios, seguidos de un segundo intento antes de S18. & Metodología; criterios 1.2, 2.1, 5.2 y 7.1, sin modificarlos \\
  Apoyo en la ficha F4 & Tarjetas de ayuda que el equipo pide cuando las necesita, sin separar a nadie de su grupo. & Metodología; criterio 7.2, sin modificarlo \\
  Ampliación y alto rendimiento & En S13, panel inclinado óptimo o más de un obstáculo en el modelo. Con los datos del socio, estimar el radio de la Tierra. & Contenidos de profundización; criterios 3.1, 6.1 \\
  TDAH & Tareas en bloques breves con hitos en Kanban, rotación ordenada para que tenga el rol de medición en la salida al patio (letra D) y más tiempo si lo necesita en el examen. & Tiempos y organización; adecuación del examen \\
  Dislexia & Enunciados en tipografía clara, lectura en voz alta, sin penalizar la ortografía en el examen y exposición oral como alternativa para la reflexión. El rol de comunicación le corresponde en el turno~1 (letra D), con textos breves y corrector. & Instrumentos y formato \\
  Incorporación reciente & Revisión de los conocimientos previos que indique el diagnóstico y compañero de referencia en el equipo. & Apoyo inicial \\
\end{tablaSAlarga}

En ningún caso se rebajan los criterios de evaluación: lo que cambian son los apoyos, los tiempos y los formatos, conforme al art.~32 del decreto.

\section{Proyecto eTwinning: «Sombras sin fronteras»}\label{sec:etwinning}

\subsection{Descripción}
eTwinning es una comunidad segura y gratuita en la que docentes y alumnado de centros europeos trabajan en proyectos colaborativos \parencite{esep}. En este proyecto, el IES Ejemplo y un centro de secundaria situado en torno a 50°~N miden en la misma semana, el mismo día si el cielo lo permite, la sombra de una varilla al mediodía solar y comparan cómo afecta la latitud a la altura del sol y a la separación de los paneles. Este planteamiento se adapta de la actividad escolar sobre Eratóstenes \parencite{eratostenes}.

\subsection{Objetivos}
\begin{itemize}
  \item Comprobar que el modelo matemático funciona en otro lugar y en otras condiciones (CE2, CE6).
  \item Comunicar datos, métodos y conclusiones a un público externo (CE8).
  \item Colaborar con otro centro de forma segura y respetuosa (CE9).
\end{itemize}

\subsection{Socio y puesta en marcha}
\begin{itemize}
  \item \textbf{Socio tipo:} grupo de edad similar en un centro de otro país eTwinning y otra latitud. La lengua de trabajo se acuerda entre ambos docentes.
  \item \textbf{Registro:} cada docente se registra con EU Login y solicita la validación de su organización nacional de apoyo \parencite{esep}.
  \item \textbf{Creación:} se necesitan dos socios fundadores; el proyecto debe ser aprobado por los Servicios Nacionales de Apoyo de ambos países \parencite{intef}. Se solicita al inicio del trimestre para que esté activo antes de S1.
\end{itemize}

\subsection{Calendario integrado en las sesiones}
\begin{tablaSA}{colspec={Q[c,wd=30pt]X[l]X[l]}, cells={font=\footnotesize}}
  Ses. & IES Ejemplo & Centro socio \\
  Antes & Registro, acuerdo de plan y calendario, apertura del TwinSpace & Igual \\
  1 & Presentación breve de cada equipo en el TwinSpace, con avatar & Presentación equivalente \\
  14 & Sombra de la varilla al mediodía solar; datos y fecha en una hoja compartida & Medición el mismo día, a su mediodía solar; con nubes, el primer día despejado de esa semana \\
  16 & Cálculo de $h$ de ambos centros; comparación de sombras y filas & Cálculo equivalente \\
  20 & Videomensaje de conclusiones y foro de preguntas & Respuesta y conclusiones conjuntas \\
\end{tablaSA}

\subsection{Resultado matemático esperado}
La altura del sol al mediodía del solsticio de invierno es $h=90^\circ-\varphi-23{,}45^\circ$. En Madrid ($40{,}4^\circ$) resulta $26{,}15^\circ$ y la sombra equivale a $2{,}04$ veces la altura del objeto; en $50^\circ$~N resulta $16{,}55^\circ$ y la sombra es $3{,}37$ veces la altura. El alumnado concluye que la misma instalación necesita más separación entre filas en el norte. El día de la medición, por ejemplo el 10 de marzo ($\delta\approx-4{,}9^\circ$), el sol está a unos $44{,}7^\circ$ en Madrid y a $35{,}1^\circ$ en $50^\circ$~N, y la varilla de 1~m da sombras de $1{,}01$~m y $1{,}42$~m. Si los centros miden en días distintos, cada altura se calcula con la declinación de su fecha. Como ampliación, con la diferencia de ángulos y la distancia entre centros sobre el mismo meridiano se estima el radio de la Tierra.

\subsection{Herramientas, seguridad y evaluación}
\begin{itemize}
  \item \textbf{Herramientas:} TwinSpace para foros y materiales, hoja de cálculo compartida, GeoGebra y videomensaje.
  \item \textbf{Seguridad y datos:} se usan avatares y nombres de pila o equipos, sin fotografías identificables del alumnado, y las autorizaciones familiares que exija el centro. El profesorado supervisa la comunicación.
  \item \textbf{Evaluación:} la comparación con el socio forma parte del libro (R1, criterios 6.1 y 2.2). El proyecto se valora con una encuesta conjunta a estudiantes y docentes sobre colaboración, aprendizaje y organización.
\end{itemize}


\printbibliography[heading=bibintoc,title={Bibliografía}]

\section*{Declaración de uso de inteligencia artificial}
\addcontentsline{toc}{section}{Declaración de uso de inteligencia artificial}
Para la elaboración de este documento se ha utilizado una herramienta de inteligencia artificial generativa (Claude, de Anthropic). La IA se ha empleado para generar el código \LaTeX{} y las matrices, buscar y contrastar fuentes, y redactar la propuesta didáctica a partir de los materiales y las decisiones del autor sobre el curso, el reto, el tipo de sesiones y la organización de la entrega.

La IA se ha utilizado únicamente como herramienta de apoyo. El alumno ha revisado, contrastado y validado todo el contenido, y asume la plena responsabilidad sobre el documento: sobre la adecuación didáctica de la propuesta, la corrección de las referencias curriculares y normativas, la idoneidad de las fuentes citadas y el cumplimiento de la normativa académica de la Universidad Rey Juan Carlos. Cualquier error u omisión que pueda contener es responsabilidad exclusiva del autor.

\campo{Autoría de la plantilla:} la plantilla del mapa ha sido proporcionada por la profesora María Mercedes Martín Lope para la asignatura Diseño Curricular de las Matemáticas del Máster en Formación del Profesorado (MAES) de la Universidad Rey Juan Carlos, curso 2026-27 \parencite{plantilla}. La propuesta sigue su guía y su rúbrica \parencite{guia,rubrica}.


\end{document}
